
\documentclass[12pt,a4paper,oneside]{report}

% ============================================================================
% PAQUETES
% ============================================================================
\usepackage[utf8]{inputenc}
\usepackage[T1]{fontenc}
\usepackage[spanish,es-nodecimaldot]{babel}
\usepackage{geometry}
\usepackage{graphicx}
\usepackage{hyperref}
\usepackage{fancyhdr}
\usepackage{tabularx}
\usepackage{longtable}
\usepackage{xcolor}
\usepackage{enumitem}
\usepackage{float}
\usepackage{booktabs}
\usepackage{multirow}
\usepackage{array}
\usepackage{colortbl}
\usepackage{caption}
\usepackage{titlesec}
\usepackage{tocloft}
\usepackage{amsmath}
\usepackage{amssymb}
\usepackage{setspace}
\usepackage{times}
\usepackage{indentfirst}
\usepackage{parskip}

% ============================================================================
% MÁRGENES (UNSAAC: izquierda 3 cm, resto 2.5 cm)
% ============================================================================
\geometry{left=3cm, right=2.5cm, top=2.5cm, bottom=2.5cm}

% ============================================================================
% INTERLINEADO
% ============================================================================
\setstretch{1.5}

% ============================================================================
% SANGRÍA Y PÁRRAFO
% ============================================================================
\setlength{\parindent}{1.25cm}
\setlength{\parskip}{6pt}

% ============================================================================
% ENCABEZADO Y PIE DE PÁGINA
% ============================================================================
\setlength{\headheight}{14pt}
\pagestyle{fancy}
\fancyhf{}
\fancyhead[R]{\small\thepage}
\fancyhead[L]{\small\textit{Sistema automatizado de transcripción de audiencias judiciales}}
\renewcommand{\headrulewidth}{0.4pt}
\fancypagestyle{plain}{
  \fancyhf{}
  \fancyhead[R]{\small\thepage}
  \renewcommand{\headrulewidth}{0.4pt}
}

% ============================================================================
% CONFIGURACIÓN DE CAPÍTULOS Y SECCIONES
% ============================================================================
\titleformat{\chapter}[display]
  {\normalfont\large\bfseries\centering}
  {CAPÍTULO \thechapter}{10pt}{\large\bfseries\centering}
\titlespacing*{\chapter}{0pt}{30pt}{20pt}

\titleformat{\section}
  {\normalfont\normalsize\bfseries}
  {\thesection}{1em}{}
\titleformat{\subsection}
  {\normalfont\normalsize\bfseries}
  {\thesubsection}{1em}{}
\titleformat{\subsubsection}
  {\normalfont\normalsize\itshape}
  {\thesubsubsection}{1em}{}

% ============================================================================
% ÍNDICE
% ============================================================================
\renewcommand{\contentsname}{ÍNDICE GENERAL}
\setcounter{tocdepth}{3}
\setcounter{secnumdepth}{3}

% ============================================================================
% TABLAS Y FIGURAS
% ============================================================================
\renewcommand{\arraystretch}{1.3}
\captionsetup{font=small, labelfont=bf, justification=raggedright, singlelinecheck=false}

% ============================================================================
% HIPERENLACES
% ============================================================================
\hypersetup{
  colorlinks=true,
  linkcolor=black,
  citecolor=black,
  urlcolor=black,
  pdftitle={Sistema automatizado de transcripción de audiencias judiciales con inteligencia artificial},
  pdfauthor={Jesús Abraham Uchiri Quispe; Carla Paloma Cueto Sánchez},
}

% ============================================================================
% INICIO DEL DOCUMENTO
% ============================================================================
\begin{document}

% ============================================================================
% PORTADA
% ============================================================================
\begin{titlepage}
\centering
\vspace*{0.5cm}

{\large\bfseries UNIVERSIDAD NACIONAL DE SAN ANTONIO ABAD DEL CUSCO\par}
\vspace{0.3cm}
{\normalsize FACULTAD DE INGENIERÍA ELÉCTRICA, ELECTRÓNICA, INFORMÁTICA Y MECÁNICA\par}
\vspace{0.2cm}
{\normalsize ESCUELA PROFESIONAL DE INGENIERÍA INFORMÁTICA Y DE SISTEMAS\par}

\vspace{1cm}

% Logo UNSAAC (incluir si se dispone del archivo)
% \includegraphics[width=4cm]{logo_unsaac.png}\\[1cm]

\vspace{0.5cm}

{\large\bfseries PLAN DE TESIS\par}
\vspace{0.8cm}

{\Large\bfseries
SISTEMA AUTOMATIZADO DE TRANSCRIPCIÓN DE AUDIENCIAS JUDICIALES\\
EMPLEANDO INTELIGENCIA ARTIFICIAL EN EL DISTRITO JUDICIAL DE CUSCO
\par}

\vspace{1.5cm}

{\normalsize
\textbf{Tesis para optar al título profesional de:}\\[0.3cm]
INGENIERO INFORMÁTICO Y DE SISTEMAS
\par}

\vspace{1cm}

{\normalsize
\textbf{Presentado por:}\\[0.3cm]
Br. JESÚS ABRAHAM UCHIRI QUISPE\\
Código: 080227\\[0.3cm]
Br. CARLA PALOMA CUETO SÁNCHEZ\\
Código: 120008
\par}

\vspace{1cm}

{\normalsize
\textbf{Asesor:}\\[0.3cm]
Lic. JOSÉ MAURO PILLCO QUISPE
\par}

\vfill

{\normalsize CUSCO -- PERÚ\\[0.2cm] \textbf{2026}}
\end{titlepage}

% ============================================================================
% PÁGINA EN BLANCO
% ============================================================================
\newpage
\thispagestyle{empty}
\mbox{}

% ============================================================================
% ÍNDICE GENERAL
% ============================================================================
\newpage
\tableofcontents

\newpage
\listoftables
\addcontentsline{toc}{chapter}{Lista de Tablas}

% ============================================================================
% INICIO DEL CONTENIDO
% ============================================================================
\newpage
\setcounter{page}{1}

% ============================================================================
% CAPÍTULO I
% ============================================================================
\chapter{PLANTEAMIENTO DE LA INVESTIGACIÓN}

\section{Planteamiento del Problema}

En el ámbito internacional, las audiencias judiciales generan grandes volúmenes de información oral que deben ser transcritos manualmente, proceso que es lento, propenso a errores humanos y que retrasa la emisión de resoluciones judiciales. Los sistemas automatizados de reconocimiento del habla (ASR) han experimentado avances significativos gracias al aprendizaje profundo, pero todavía no alcanzan la precisión humana en condiciones reales. Estudios recientes muestran que incluso los modelos más robustos, como Whisper y Deepgram Nova, mantienen tasas de error significativas especialmente en escenarios con ruido de fondo, acentos variables y habla espontánea, mientras que la transcripción humana profesional supera consistentemente el 98\,\% de exactitud (Radford et al., 2023). Esta brecha subsiste especialmente en entornos con ruido de fondo y diversidad de acentos, donde la supervisión humana sigue siendo indispensable para obtener registros fieles, coherentes y utilizables en contextos sensibles como los ámbitos jurídico, clínico y administrativo.

Evidencias empíricas demuestran que los principales sistemas ASR comerciales presentan brechas significativas de desempeño entre grupos poblacionales. El estudio de Koenecke et al. (2020), basado en entrevistas sociolingüísticas con 73 hablantes afrodescendientes y 42 hablantes blancos, revela que la tasa de error de palabras (WER) promedio asciende a \textbf{0,35} para hablantes afrodescendientes, mientras que se reduce a \textbf{0,19} para hablantes blancos, incluso cuando se analizan fragmentos de audio comparables en edad, género y duración. Esta disparidad refleja sesgos estructurales en los modelos acústicos y afecta negativamente su idoneidad para los contextos judiciales donde la precisión es fundamental.

En América Latina, la adopción de sistemas de reconocimiento automático del habla en entornos judiciales es aún incipiente, debido a la falta de infraestructura tecnológica. Estudios sobre digitalización del sector público en regiones en desarrollo, como el análisis de Bassey, Mulligan y Ojo (2022), evidencian que los países latinoamericanos presentan brechas profundas en infraestructura digital, capacidades institucionales y preparación tecnológica, factores que dificultan el despliegue de herramientas basadas en inteligencia artificial en procesos críticos de gestión documental y registro oral. Además, aunque la región representa aproximadamente el 14\,\% de las visitas globales a soluciones de inteligencia artificial, solo recibe alrededor del 1,12\,\% de la inversión mundial en IA, lo que evidencia una brecha de preparación y despliegue tecnológico (CEPAL, 2025).

En el Perú, y específicamente en el Distrito Judicial de Cusco, las audiencias se registran mediante transcripción manual, lo que genera retrasos, errores y carga administrativa sobre el personal judicial. Hasta la fecha, no existen estudios académicos que evalúen sistemas ASR aplicados al sector judicial peruano. El estudio de Zuchowski y Göller (2022), basado en la observación directa de 313 muestras de documentación clínica, encontró que la transcripción mediante ASR permite completar formularios en un promedio de 5,11 minutos, frente a los 8,9 minutos requeridos por la digitación tradicional, lo que representa una mejora del 43\,\% en eficiencia temporal. Estos hallazgos sugieren que la implementación de ASR en los órganos jurisdiccionales del Cusco podría reducir significativamente los tiempos de procesamiento documental, mejorar la calidad del registro de audiencias y aliviar la sobrecarga que actualmente enfrenta el personal judicial.

En el Distrito Judicial de Cusco, los trabajadores judiciales —jueces, secretarios y especialistas de audiencias— enfrentan dificultades derivadas de la transcripción manual de las audiencias. Esta práctica, que aún constituye la norma en los procesos judiciales peruanos, ocasiona errores frecuentes, retrasos en la redacción de actas de audiencia y una sobrecarga significativa en el personal encargado de dicha labor. Como resultado, se pone en riesgo la calidad y confiabilidad de los registros procesales, situación que adquiere especial relevancia actualmente, cuando se intensifica el uso de audiencias virtuales como parte de la modernización de la justicia.

Las causas de esta problemática son múltiples. En primer lugar, persiste una fuerte dependencia de métodos manuales tradicionales que ralentizan la labor judicial. A ello se suma la falta de implementación de sistemas automatizados con inteligencia artificial en el Poder Judicial, específicamente en el Distrito Judicial de Cusco, lo que limita el aprovechamiento de las tendencias tecnológicas aplicadas al sector justicia. Asimismo, la capacitación tecnológica del personal judicial sigue siendo insuficiente, dificultando la adaptación a nuevas herramientas digitales. Finalmente, el incremento en el volumen de audiencias virtuales tras la pandemia ha generado una saturación administrativa que agrava la situación existente.

Las consecuencias de mantener este escenario son evidentes. El retraso en la emisión de actas y resoluciones judiciales atenta contra el principio de celeridad procesal, mientras que los errores en la transcripción pueden distorsionar la información, comprometiendo el derecho al debido proceso. La sobrecarga laboral del personal judicial incrementa el riesgo de desgaste profesional y reduce la eficiencia institucional. En un nivel más amplio, esta situación alimenta la percepción de ineficiencia y opacidad, generando desconfianza ciudadana respecto a la transparencia del sistema de justicia.

Además, el lanzamiento en 2023 del aplicativo \textit{Transcriptor de Audiencias Versión Alfa 2} por parte del Poder Judicial del Perú, implementado inicialmente en la Corte Superior de Justicia de Sullana y, posteriormente, en Huaura, buscó automatizar la transcripción de audiencias virtuales y reducir la carga procesal del personal judicial, constituyendo una experiencia práctica pionera en el país (Poder Judicial del Perú, 2023). Sin embargo, a la fecha no se dispone de publicaciones académicas ni reportes técnicos que documenten su precisión, funcionamiento lingüístico o adaptabilidad.

Frente a este panorama, la presente investigación no se limita a evaluar una herramienta preexistente de transcripción automática, sino que propone el diseño y desarrollo de un prototipo funcional de un sistema automatizado de transcripción judicial basado en inteligencia artificial, que integra dos servicios de IA: Deepgram Nova-3 para el reconocimiento automático del habla en tiempo real con detección de múltiples hablantes, y Claude Sonnet 4 (Anthropic) para el mejoramiento lingüístico del texto y la generación de actas con formato oficial. Complementariamente, se emplea faster-whisper (variante optimizada de Whisper de OpenAI) para el reprocesamiento por lotes de audios pregrabados. El sistema se valida mediante un corpus lingüístico judicial construido a partir de grabaciones reales del Distrito Judicial de Cusco, permitiendo evaluar la precisión frente a variaciones de acento, ruido y terminología jurídica. De este modo, el aporte principal radica en la integración y evaluación empírica de un sistema de IA contextualizado al ámbito judicial, diferenciándose de los estudios previos centrados únicamente en la comparación de herramientas comerciales.

\section{Formulación de la Pregunta de Investigación}

\subsection{Problema General}

¿En qué medida un sistema automatizado de transcripción de audiencias judiciales empleando inteligencia artificial puede mejorar la precisión, la calidad y la eficiencia del proceso de transcripción en el Distrito Judicial de Cusco?

\subsection{Problemas Específicos}

\begin{enumerate}[label=\alph*)]
  \item ¿Cuáles son las limitaciones y dificultades del proceso actual de transcripción manual de audiencias judiciales en el Distrito Judicial de Cusco?
  \item ¿Qué características técnicas y funcionales debe poseer un sistema automatizado de transcripción judicial basado en inteligencia artificial para garantizar precisión y eficiencia?
  \item ¿Cómo se puede diseñar e implementar un sistema automatizado de transcripción adaptado a las condiciones tecnológicas y operativas del Distrito Judicial de Cusco?
  \item ¿En qué medida la implementación del sistema automatizado mejora la precisión, reduce los tiempos de transcripción y eleva la satisfacción de los usuarios en comparación con el método manual?
\end{enumerate}

\section{Objetivos de la Investigación}

\subsection{Objetivo General}

Diseñar, implementar y evaluar un sistema automatizado de transcripción de audiencias judiciales empleando inteligencia artificial, que permita mejorar la precisión, calidad y eficiencia del proceso de transcripción en el Distrito Judicial de Cusco.

\subsection{Objetivos Específicos}

\begin{enumerate}[label=\alph*)]
  \item Analizar las limitaciones y dificultades del proceso actual de transcripción manual de audiencias judiciales en el Distrito Judicial de Cusco, a fin de identificar los principales factores que afectan la precisión, calidad y eficiencia del registro judicial.
  \item Diseñar la arquitectura técnica y funcional de un sistema automatizado de transcripción judicial basado en inteligencia artificial, estableciendo sus módulos, flujos de datos, requerimientos de hardware y software, y niveles de precisión esperados.
  \item Implementar el prototipo del sistema automatizado de transcripción judicial empleando inteligencia artificial, integrando componentes de captura de audio, procesamiento, comunicación cliente-servidor y motor de reconocimiento de voz.
  \item Evaluar el desempeño del sistema automatizado en comparación con el método manual, utilizando métricas cuantitativas de precisión y rendimiento.
\end{enumerate}

\section{Justificación de la Investigación}

\subsection{Por su Conveniencia}

El presente estudio se justifica por la utilidad y las aplicaciones de la inteligencia artificial en el ámbito jurídico y por la necesidad urgente de mejorar el tiempo y la precisión en el registro de las actas de audiencias. La transcripción manual de audiencias en el Distrito Judicial de Cusco ha demostrado ser lenta, propensa a errores y con una carga administrativa excesiva que afecta directamente la continuidad procesal y la calidad del servicio de justicia. Implementar un sistema automatizado basado en ASR resulta conveniente porque permitirá superar estas limitaciones y alinear el Poder Judicial con tendencias globales de digitalización. En experiencias internacionales, como en el Parlamento de Portugal, la adopción de sistemas de transcripción automática redujo significativamente los errores y mejoró la transparencia en los procesos legislativos (Nascimento et al., 2024).

Bajo la perspectiva metodológica, la conveniencia de este trabajo también radica en su factibilidad y pertinencia, mientras que resulta necesario porque responde a una demanda institucional concreta que afecta tanto a operadores de justicia como a los usuarios del sistema.

\subsection{Por su Relevancia}

La relevancia de la presente investigación reside en el impacto directo que tendrá sobre la calidad y precisión de las transcripciones de las audiencias, aspecto vital para garantizar el derecho al debido proceso. En América Latina y en el Perú, aún no se han desarrollado estudios empíricos suficientes que evalúen la eficacia de los sistemas ASR en contextos judiciales, lo que revela un vacío académico que esta investigación busca abordar. Koenecke et al. (2020), desde \textit{PNAS}, evidencian que estos sistemas presentan disparidades de precisión según características sociolingüísticas, lo que subraya la necesidad de adaptarlos y evaluarlos en contextos específicos, como el judicial peruano.

Desde la visión de esta investigación, la relevancia se fortalece al evidenciar cómo este trabajo contribuye a solucionar una problemática real con implicancias sociales y jurídicas; además, es relevante porque permite reforzar la confianza ciudadana en la justicia, optimizando procesos internos y beneficiando a todos los actores del sistema judicial.

\subsection{Por sus Implicancias Prácticas}

Las implicancias prácticas de esta investigación son múltiples. A corto plazo, permitirá obtener actas judiciales con mayor rapidez y exactitud, reduciendo la sobrecarga del personal judicial. A mediano y largo plazo, el modelo puede replicarse en otros distritos judiciales del país, generando un estándar tecnológico de justicia digital. Legchenko y Bondarenko (2025) comprobaron que entrenar modelos ASR con corpus lingüísticos especializados incrementa la precisión de las transcripciones en contextos jurídicos, lo que abre la posibilidad de integrar, en el futuro, sistemas de análisis semántico y búsqueda automática de jurisprudencia.

Estas implicancias se justifican por el aporte a la transformación digital del Poder Judicial y por el beneficio directo que otorgan a jueces, secretarios, fiscales y ciudadanos que requieren procesos más céleres y confiables.

\subsection{Por su Valor Teórico}

En el plano teórico, este estudio aporta al desarrollo del conocimiento aplicado al derecho, al demostrar cómo la inteligencia artificial puede incidir en la calidad de las transcripciones judiciales. Aunque existen investigaciones en parlamentos europeos y en cortes supremas anglosajonas (Legchenko y Bondarenko, 2025), son escasos los estudios que exploran su aplicación en el ámbito latinoamericano y, particularmente, en el Perú. Esto convierte a la investigación en un referente pionero.

El valor teórico se centra en demostrar la relación entre variables —la implementación de IA y la calidad de las transcripciones—, además radica en vincular la innovación tecnológica con la garantía de derechos fundamentales y el fortalecimiento de la justicia digital.

\subsection{Por su Valor Metodológico}

Finalmente, la investigación tiene una utilidad metodológica relevante, pues adopta un diseño preexperimental de un grupo, lo que permitirá medir objetivamente el impacto del sistema de IA en la calidad de las transcripciones judiciales. Esta metodología puede ser reutilizada en otros estudios, sirviendo como referencia para futuras investigaciones en contextos judiciales o administrativos.

La utilidad metodológica radica en su capacidad de ser replicada y validada en distintos escenarios; se justifica porque permite garantizar la transparencia de los resultados y asegurar que la innovación propuesta pueda sostenerse en investigaciones futuras.

\section{Delimitación del Estudio}

\subsection{Delimitación Temporal}

La delimitación temporal de la presente investigación comprende el periodo entre noviembre de 2025 y junio de 2026, que coincide con la fase de auge de las audiencias virtuales y semipresenciales posteriores a la pandemia, etapa en la que se evidenció una mayor carga operativa y dificultades en la gestión documental judicial. Asimismo, se constituye en el marco temporal en el cual se recopilan los datos de audiencias, se realizan las observaciones institucionales, se implementa el prototipo del sistema y se evalúa la factibilidad técnica de aplicar sistemas de reconocimiento automático del habla en los módulos judiciales.

\subsection{Delimitación Espacial}

El presente estudio se llevará a cabo en el Distrito Judicial de Cusco, Perú, específicamente en las salas donde se desarrollan audiencias. Este distrito judicial ha incrementado de manera significativa el uso de plataformas digitales en la administración de justicia tras la pandemia de COVID-19, lo que ha puesto en evidencia limitaciones en los métodos manuales de transcripción y la necesidad de adoptar tecnologías emergentes.

La elección del Distrito Judicial de Cusco se justifica porque representa un contexto donde la modernización digital es urgente y se enfrenta a un volumen creciente de audiencias.

% ============================================================================
% CAPÍTULO II
% ============================================================================
\chapter{MARCO TEÓRICO}

\section{Antecedentes}

\subsection{Antecedentes a Nivel Internacional}

Garneau y Bolduc (2024), en su estudio titulado \textit{The State of Commercial Automatic French Legal Speech Recognition Systems and their Impact on Court Reporters}, realizado en la University of Copenhagen, se centraron en analizar el impacto de los sistemas de reconocimiento automático del habla (ASR) en el ámbito judicial canadiense y, particularmente, en el francés legal. El objetivo general fue evaluar la viabilidad de emplear ASR en la transcripción de procesos judiciales, considerando tanto su rendimiento técnico como las implicancias profesionales y legales derivadas de su implementación.

El diseño de la investigación fue experimental comparativo, basado en la evaluación de tres modelos de ASR (Amazon Transcribe, Google Speech-to-Text y OpenAI Whisper). La muestra se conformó por audios reales de carácter judicial, incluyendo grabaciones de la \textit{Commission Charbonneau} y sentencias dictadas por jueces, con un total aproximado de 1 día, 6 horas y 37 minutos de material de audio. Como instrumentos, se utilizaron transcripciones oficiales realizadas por taquígrafos profesionales, que sirvieron como referencia para contrastar los resultados de los sistemas automáticos.

La investigación utilizó principalmente la métrica de Word Error Rate (WER) para evaluar la precisión de los modelos, considerando inserciones, eliminaciones y sustituciones. Asimismo, introdujo la Sonnex Distance, una medida fonética que permite identificar errores relacionados con homófonos, género, número y variaciones dialectales, complementando así el análisis más allá del conteo léxico.

Los hallazgos demostraron que AWS obtuvo el mejor desempeño promedio con un WER de 0,15, seguido de Google y Whisper. Sin embargo, ninguno de los sistemas alcanzó la precisión requerida para prescindir de supervisión humana en el ámbito judicial. En conclusión, los autores destacan que, aunque prometedores, los sistemas ASR aún deben especializarse y perfeccionarse para satisfacer las demandas críticas de exactitud en el contexto legal.

Harrington (2023), en su artículo titulado \textit{Incorporating automatic speech recognition methods into the transcription of police-suspect interviews: factors affecting automatic performance}, publicado en la Universidad de York (Reino Unido), aborda la problemática de la baja calidad y subjetividad en las transcripciones ortográficas de entrevistas policiales entre sospechosos y autoridades judiciales. El objetivo general fue evaluar la viabilidad de integrar métodos de reconocimiento automático de voz (ASR) en el proceso de transcripción de entrevistas policiales, con la finalidad de generar registros textuales más fiables y eficientes para su utilización como evidencia en tribunales.

El diseño de la investigación fue de tipo experimental comparativo. La muestra estuvo compuesta por grabaciones extraídas de dos corpus de inglés británico forensemente relevantes: el \textit{Dynamic Variability in Speech database} (DyViS), con 100 hablantes del sur de Inglaterra, y el \textit{West Yorkshire Regional English Database} (WYRED), con 180 hablantes del norte de Inglaterra. Se utilizaron tres sistemas de reconocimiento automático de voz: Amazon Transcribe, Google Cloud Speech-to-Text y Rev AI.

Los resultados mostraron que Amazon Transcribe obtuvo el mejor desempeño global, con tasas de error menores al 14\,\% en condiciones de audio limpio, mientras que Google presentó el peor rendimiento. La autora concluye que los sistemas automáticos podrían utilizarse como un primer borrador de transcripciones, siempre que exista un proceso posterior de edición humana.

Chaloupka, Palecek y Cerva (2021), en su investigación titulada \textit{Audio-visual broadcast transcription system using artificial neural networks}, publicada en el IEEE International Workshop on Electronics, Control, Measurement, Signals and their Application to Mechatronics, presentaron un sistema innovador para la transcripción de noticias televisivas utilizando redes neuronales artificiales. El objetivo general fue diseñar y evaluar un modelo de transcripción automática que integrara simultáneamente el procesamiento de audio y la extracción de información visual, con el fin de mejorar la precisión de las transcripciones en entornos complejos.

La muestra utilizada consistió en un corpus audiovisual de 817 horas de grabaciones de noticieros de cinco canales de televisión checos. Para la evaluación, se aplicaron métricas estándar como el Word Error Rate (WER), reduciéndolo de un 29,4\,\% inicial a 15,4\,\% tras aplicar algoritmos de corrección léxica y técnicas de preprocesamiento visual.

Los resultados evidenciaron que el sistema híbrido logró una transcripción más precisa y robusta en comparación con sistemas de audio únicamente. Los autores concluyen que este tipo de sistemas tiene un alto potencial para ser aplicado en contextos judiciales, mediáticos y administrativos.

Loakes (2024), en su artículo titulado \textit{Automatic speech recognition and the transcription of indistinct forensic audio: how do the new generation of systems fare?}, publicado en \textit{Frontiers in Communication}, llevó a cabo un análisis actualizado sobre el desempeño de los sistemas ASR en contextos de audio forense indistinto. Los resultados mostraron que Whisper obtuvo el mejor rendimiento en condiciones adversas, logrando transcribir el 68,9\,\% del material, con un WER del 50\,\%. Las conclusiones indicaron que la transcripción humana continúa siendo indispensable hasta que los sistemas automáticos alcancen un grado de confiabilidad superior en condiciones adversas de audio.

\subsection{Antecedentes a Nivel Nacional}

Vera (2021) desarrolló un estudio titulado \textit{STT: Un sistema de apoyo a la transcripción de audiencias fiscales usando Vosk}, publicado en la \textit{Revista de Investigación de Sistemas e Informática} de la Universidad Nacional Mayor de San Marcos. El objetivo general consistió en diseñar e implementar un modelo transcriptor en idioma español que permitiera agilizar el proceso de transcripción de audiencias fiscales. El desarrollo del sistema se apoyó en una arquitectura cliente-servidor utilizando VueJS para la interfaz, FastAPI para el backend, RabbitMQ y Celery como gestores de tareas asíncronas, y PostgreSQL para la gestión de datos. Los resultados mostraron que Vosk fue elegido como motor principal del sistema por su eficiencia y rapidez.

Larrea (2024), en su tesis de maestría titulada \textit{Uso de la plataforma Google Meet y transcriptor IA en el proceso de trámite de las audiencias del módulo penal dentro del Poder Judicial de Junín en la sub sede central Huancayo 2024}, tuvo como objetivo general evaluar la influencia de la plataforma Google Meet y un transcriptor basado en inteligencia artificial en la gestión de audiencias penales. El diseño fue cuasiexperimental de grupo único con mediciones pre y post implementación. Los resultados evidenciaron que la implementación contribuyó a una mayor celeridad en las audiencias virtuales frente a las presenciales.

Rubio Portilla (2023), en su artículo \textit{Audiencias virtuales y su impacto en el debido proceso penal en un distrito judicial}, publicado en la revista \textit{SCIÉNDO}, determinó el impacto de las audiencias virtuales en la garantía del debido proceso penal. Se aplicó la prueba estadística de Pearson, obteniendo una correlación significativa entre el uso de audiencias virtuales y el respeto del debido proceso penal (r = 0,837; p < 0,05).

En 2023, el Poder Judicial del Perú lanzó el aplicativo \textit{Transcriptor de Audiencias Versión Alfa 2} en la Corte Superior de Justicia de Sullana, con el objetivo de disminuir la carga procesal mediante la generación casi inmediata de transcripciones al concluir las audiencias virtuales. El uso de la herramienta disminuyó el tiempo de elaboración en un 30\,\% y redujo errores manuales en un 15\,\% (Poder Judicial del Perú, 2023).

\subsection{Antecedentes a Nivel Local}

Yépez-Reyes y Cruz-Silva (2024), en el artículo titulado \textit{Inteligencia artificial en la transcripción de entrevistas}, publicado en la revista \textit{Contratexto}, plantearon como objetivo principal identificar la herramienta de inteligencia artificial más adecuada para la transcripción de entrevistas grabadas en español, priorizando la completitud de tareas, la eficiencia y la eficacia.

Los investigadores analizaron cuatro plataformas de IA: Office 365 (Word) Transcribe, Amazon Transcribe, Notta y Whisper (OpenAI). Los resultados demostraron que, aunque las plataformas de IA permiten transcripciones rápidas y precisas, cada una presenta ventajas y limitaciones específicas. Se concluyó que la transcripción mediante IA no sustituye completamente la revisión humana, pero sí permite optimizar tiempos y mejorar la productividad.

\section{Marco Normativo y Operativo del Proceso de Audiencia Judicial}

\subsection{Concepto de Sistemas Automatizados en Entornos Judiciales}

Los sistemas automatizados en entornos judiciales se entienden como aplicaciones tecnológicas basadas en inteligencia artificial y algoritmos de procesamiento de información, cuyo propósito es apoyar o sustituir tareas rutinarias realizadas por los operadores de justicia. Estos sistemas abarcan desde la transcripción automática de audiencias hasta la clasificación documental y la gestión de expedientes digitales. Según Cabrera Fernández (2024), la automatización judicial constituye un medio para optimizar la celeridad procesal y fortalecer la transparencia institucional, siempre que se mantenga la supervisión humana en la toma de decisiones complejas.

La incorporación de sistemas ASR en tribunales constituye un avance significativo hacia la justicia digital, aunque aún persisten limitaciones de precisión y sesgos en contextos multilingües. Loakes (2024) advierte que los sistemas automatizados enfrentan problemas de exactitud en audios de baja calidad, lo que condiciona su fiabilidad como herramientas únicas para generar registros oficiales.

\subsection{Valor Legal de las Actas de Audiencia}

Las actas de audiencia constituyen documentos públicos de carácter procesal y tienen plena validez legal en el marco de los procesos judiciales peruanos. Según el Código Procesal Civil (art. 197) y el Código Procesal Penal (art. 120), las actas constituyen prueba preconstituida y registran de manera fidedigna los actos procesales, garantizando el principio de publicidad y el derecho de defensa.

El Poder Judicial del Perú (2023) ha señalado que la digitalización de las actas, a través de grabaciones audiovisuales y su posterior transcripción, debe cumplir con estándares de autenticidad, integridad y disponibilidad, tal como exige la Ley de Gobierno Digital N.\textordmasculine{} 1412 y el Decreto Legislativo N.\textordmasculine{} 681 sobre microformas digitales con valor legal.

\subsection{Normativa Aplicable al Uso de Sistemas Automatizados}

El marco normativo vigente en el Perú ha comenzado a reconocer el uso de medios electrónicos y tecnologías emergentes en la justicia. El Reglamento de Audiencias Virtuales (R.A. 000-2020-CE-PJ) y las disposiciones del Consejo Ejecutivo del Poder Judicial regulan la conducción de audiencias virtuales, el registro audiovisual y la protección de datos personales de las partes involucradas.

Para la transcripción automatizada, se debe considerar la Ley N.\textordmasculine{} 29733 --Ley de Protección de Datos Personales--, que obliga a garantizar la privacidad de los intervinientes; la Ley de Firmas y Certificados Digitales (D.S. N.\textordmasculine{} 052-2008-PCM), que respalda el uso de firmas digitales en documentos generados electrónicamente; y la NTP 392.030-2:2015 (INACAL), que regula la producción y almacenamiento de microformas digitales con valor legal.

\subsection{Procedimiento en una Sala de Audiencia}

El desarrollo de una audiencia judicial sigue etapas formales que deben ser registradas en su integridad: instalación de la audiencia, lectura de cargos y alegatos de apertura, práctica de pruebas, alegatos de clausura y decisión del juez. Durante todo el proceso, el secretario judicial o especialista de audiencia es responsable de la redacción del acta, asegurando que los hechos y decisiones queden reflejados sin distorsiones.

\subsection{Adaptación del Sistema ASR para Uso Judicial}

La adaptación de los motores ASR al dominio judicial no requiere el reentrenamiento de los modelos base, dado que el sistema emplea servicios preentrenados (Deepgram Nova-3 y faster-whisper large-v3) que ya ofrecen alta precisión en español. En su lugar, la adaptación se realiza mediante tres estrategias complementarias: (a) configuración de \textit{keyterms} jurídicos, enviando hasta 100 términos especializados a Deepgram en cada sesión de transcripción para mejorar el reconocimiento de vocabulario procesal (por ejemplo, sobreseimiento, requerimiento acusatorio, Señoría); (b) implementación de un diccionario jurídico local con búsqueda fuzzy (Soundex adaptado al español y distancia de Levenshtein) para la detección y corrección post-transcripción de variantes erróneas; y (c) mejoramiento lingüístico mediante Claude Sonnet 4, que corrige puntuación, aplica capitalización a términos jurídicos y reemplaza palabras mal transcritas sin añadir contenido.

La evaluación del desempeño se realiza mediante métricas estándar: WER para la precisión léxica, Sonnex Distance para la similitud fonética, Real-Time Factor (RTF) para la latencia, y MOS para la calidad percibida por los operadores judiciales.

\subsection{Aspectos de Redacción y Estilo en Actas}

El lenguaje en las actas debe ser preciso, impersonal y formal. La Guía de Redacción Judicial (CNM, 2022) recomienda el uso de tercera persona y tiempo presente, estructura clara y numerada por actos procesales, e inclusión de firmas digitales o electrónicas para cerrar el documento.

\subsection{Consideraciones Éticas y de Transparencia}

El uso de IA en justicia debe regirse por principios de equidad, explicabilidad y supervisión humana (OCDE, 2023). El juez y las partes deben ser informados de que se está utilizando un sistema automatizado para transcripción y deben tener derecho a impugnar el acta si se identifican errores.

\subsection{Inteligencia Artificial Aplicada al Reconocimiento de Voz}

En los últimos años, el reconocimiento de voz ha experimentado avances significativos tanto en el ámbito académico como en el comercial. En el plano académico, Radford et al. (2022) desarrollaron el modelo \textit{Whisper}, entrenado con aproximadamente 680\,000 horas de datos multilingües y multitarea, lo cual permitió obtener un sistema robusto frente a variaciones de acento, ruido de fondo y terminología técnica. Sin embargo, Whisper opera exclusivamente en modo offline (por lotes), lo que limita su aplicabilidad en escenarios que requieren transcripción en tiempo real.

En el ámbito comercial, servicios como \textbf{Deepgram Nova-3} han abordado esta limitación ofreciendo transcripción en tiempo real (\textit{streaming}) mediante conexión WebSocket, con latencias inferiores a 300\,ms y funcionalidades nativas de diarización de hablantes, indicadores de confianza por palabra y soporte para términos personalizados (\textit{keyterms}). Estas características son fundamentales para el ámbito judicial, donde la transcripción debe reflejar las intervenciones de múltiples participantes en tiempo real y con precisión terminológica. La elección de Deepgram Nova-3 como motor principal de transcripción en el sistema propuesto se justifica por tres razones: (1) capacidad de streaming nativo con baja latencia, indispensable para audiencias en vivo; (2) diarización integrada que identifica automáticamente a los hablantes sin componentes adicionales; y (3) soporte de keyterms que permite adaptar el reconocimiento al vocabulario jurídico sin reentrenamiento del modelo.

Complementariamente, el modelo Whisper se emplea en su variante optimizada \textbf{faster-whisper} (CTranslate2) exclusivamente para el procesamiento por lotes de audios pregrabados, donde la mayor precisión en condiciones controladas compensa la ausencia de capacidad de streaming.

\subsection{Calidad en la Transcripción de Audiencias Virtuales}

La calidad de las transcripciones en las audiencias judiciales virtuales se ha convertido en un eje crítico para la transparencia y eficiencia del sistema de justicia. Lyra et al. (2024) desarrollaron un sistema automatizado de transcripción que alcanzó un nivel de acierto cercano al 92\,\% en audios de buena calidad, lo que permitió disminuir significativamente el esfuerzo administrativo del personal.

Por su parte, Bild et al. (2021) subrayan la relevancia del componente acústico en audiencias virtuales, mostrando que una baja calidad del audio no solo afecta la fidelidad de la transcripción, sino también la percepción de credibilidad hacia los testigos.

\section{Diseños de Sistemas de Transcripción Automática}

Los avances en el diseño de sistemas de transcripción automática de audios judiciales han girado en torno a arquitecturas híbridas que combinan redes neuronales profundas, modelos ocultos de Markov y técnicas de preprocesamiento de señal.

\subsection{Modelos de Diseño Híbrido DNN-HMM}

El modelo híbrido \textit{Deep Neural Networks -- Hidden Markov Models} (DNN-HMM) ha demostrado ser una de las arquitecturas más eficaces en el campo del reconocimiento automático del habla. Este diseño integra la capacidad de representación no lineal de las redes neuronales profundas con la solidez estadística de los modelos ocultos de Markov, lo que permite capturar tanto las dependencias acústicas complejas como las transiciones temporales del habla. Chaloupka, Palecek y Cerva (2021) lo aplicaron en un sistema de transcripción audiovisual para noticias televisivas, alcanzando un desempeño superior en la reducción de errores de palabra en comparación con sistemas basados únicamente en modelos acústicos tradicionales.

Desde el punto de vista arquitectónico, los autores implementaron una red neuronal profunda con cinco capas ocultas (1024-1024-768-768-512 neuronas) y funciones de activación ReLU, acompañada de una capa de salida con función Softmax para la clasificación de estados fonéticos. El entrenamiento se realizó en 50 épocas, integrando un vocabulario de más de 500\,000 palabras y un modelo de lenguaje n-gram con suavizado Kneser-Ney.

En cuanto a métricas, la principal medida empleada fue la Word Error Rate (WER), que cuantifica errores de inserción, eliminación y sustitución. Inicialmente, los sistemas basados en DNN-HMM reportaban un WER del 29,4\,\%, pero mediante la incorporación de información visual y estrategias de preprocesamiento los investigadores lograron reducirlo hasta un 15,4\,\%.

\begin{table}[H]
\centering
\caption{Arquitectura del modelo híbrido DNN-HMM}
\begin{tabular}{|p{3.5cm}|p{4cm}|p{4cm}|}
\hline
\rowcolor{gray!20}
\textbf{Componente} & \textbf{Descripción} & \textbf{Detalles técnicos} \\
\hline
Entradas & Características acústicas & MFCC, 44.1 kHz, 16 bits \\
\hline
Red DNN & 5 capas ocultas & 1024-1024-768-768-512 neuronas, ReLU \\
\hline
Capa de salida & Clasificación de estados fonéticos & Softmax \\
\hline
Modelo de lenguaje & N-gram & Suavizado Kneser-Ney, 500k palabras \\
\hline
Entrenamiento & Iteraciones & 50 épocas \\
\hline
Métricas & Precisión & WER reducido de 29,4\,\% a 15,4\,\% \\
\hline
\end{tabular}
\label{tab:dnn_hmm}
\end{table}

\subsection{Modelos ASR Multilingües y Servicios Comerciales}

El campo del reconocimiento automático del habla ha sido transformado tanto por modelos académicos de código abierto como por servicios comerciales basados en la nube. Entre los modelos académicos, Radford et al. (2022) destacan que Whisper logra generalizar de manera efectiva en condiciones no vistas durante el entrenamiento, superando a modelos DNN-HMM tradicionales. En el ámbito comercial, servicios como Deepgram Nova-3 ofrecen transcripción en tiempo real (streaming) con latencias inferiores a 300\,ms y funcionalidades nativas de diarización, lo que los hace especialmente adecuados para aplicaciones judiciales en vivo donde la inmediatez es crítica.

\begin{table}[H]
\centering
\caption{Comparación de modelos ASR mulitlingües y comerciales}
\begin{tabular}{|p{2.5cm}|p{4cm}|p{2.5cm}|p{3.5cm}|}
\hline
\rowcolor{gray!20}
\textbf{Modelo} & \textbf{Diseño / Arquitectura} & \textbf{Idiomas Soportados} & \textbf{Características Relevantes} \\
\hline
Deepgram Nova-3 & Modelo propietario basado en deep learning, optimizado para streaming en tiempo real & +30 idiomas (incl. es-419) & Diarización nativa, keyterms, confianza por palabra, baja latencia ($< 300$ ms) \\
\hline
Whisper (OpenAI) & Codificador-Decodificador Transformer con autoatención multinivel & +90 idiomas & Entrenamiento multitarea (traducción + transcripción); robustez en ruido y acentos \\
\hline
Wav2Vec 2.0 (Meta) & Modelo auto-supervisado basado en Transformers & 53 idiomas & Preentrenamiento sin etiquetas + ajuste fino en corpus de voz etiquetado \\
\hline
SpeechT5 (Microsoft) & Transformer multimodal con soporte para síntesis y reconocimiento & Multilingüe & Capacidad de convertir texto a voz y viceversa en un mismo modelo \\
\hline
\end{tabular}
\label{tab:transformers}
\end{table}

\begin{table}[H]
\centering
\caption{Métricas de desempeño de modelos ASR}
\begin{tabular}{|p{2.5cm}|p{3.5cm}|p{3.5cm}|p{2.5cm}|}
\hline
\rowcolor{gray!20}
\textbf{Modelo} & \textbf{WER en Audio Limpio} & \textbf{WER en Audio con Ruido} & \textbf{Latencia Promedio} \\
\hline
Deepgram Nova-3 & 8,4\,\% (español, benchmark interno) & $\sim$20\,\% en ruido moderado & $< 300$ ms en streaming \\
\hline
Whisper (Large) & 2,7\,\% (LibriSpeech) & 50\,\% (audio forense, Loakes, 2024) & Solo offline (no streaming nativo) \\
\hline
Wav2Vec 2.0 Base & 5,2\,\% (LibriSpeech) & 58\,\% en ruido simulado & $<$ 2 s en inferencia offline \\
\hline
SpeechT5 & 4,8\,\% (VoxPopuli) & 55\,\% en ruido moderado & 2,5 s promedio \\
\hline
\end{tabular}
\label{tab:metricas_transformers}
\end{table}

\begin{table}[H]
\centering
\caption{Parámetros y estrategias de entrenamiento}
\begin{tabular}{|p{2.5cm}|p{3cm}|p{3.5cm}|p{3.5cm}|}
\hline
\rowcolor{gray!20}
\textbf{Modelo} & \textbf{Tamaño de Parámetros} & \textbf{Datos de Entrenamiento} & \textbf{Estrategias de Optimización} \\
\hline
Whisper & 1,5B -- 1,55B & 680\,000 h de audio multilingüe & Aprendizaje multitarea, \textit{data augmentation} con ruido \\
\hline
Wav2Vec 2.0 & 95M -- 317M & 53 idiomas (VoxPopuli) & Preentrenamiento auto-supervisado, \textit{fine-tuning} supervisado \\
\hline
SpeechT5 & 280M & Multimodal (texto + voz) & \textit{Transfer learning} y ajuste fino en tareas específicas \\
\hline
\end{tabular}
\label{tab:entrenamiento}
\end{table}

\begin{table}[H]
\centering
\caption{Validación y tipos de error evaluados}
\begin{tabular}{|p{2.5cm}|p{3cm}|p{3.5cm}|p{3.5cm}|}
\hline
\rowcolor{gray!20}
\textbf{Modelo} & \textbf{Tipos de Error Medidos} & \textbf{Herramientas de Evaluación} & \textbf{Resultados Clave} \\
\hline
Deepgram Nova-3 & Inserciones, sustituciones, omisiones, errores de diarización & WER, confianza por palabra & Alta precisión en streaming; diarización nativa con identificación de hablante \\
\hline
Whisper & Inserciones, sustituciones, omisiones & WER, BLEU para traducción & Alta precisión en entornos controlados, disminuye en audio forense \\
\hline
Wav2Vec 2.0 & Errores fonéticos y semánticos & WER, Character Error Rate (CER) & Mejor rendimiento en entornos de bajo ruido \\
\hline
SpeechT5 & Errores de prosodia y segmentación & MOS (Mean Opinion Score), WER & Buen balance entre naturalidad y exactitud \\
\hline
\end{tabular}
\label{tab:validacion}
\end{table}

\subsection{Modelos Cliente-Servidor y Arquitecturas Distribuidas}

Los modelos cliente-servidor constituyen uno de los pilares fundamentales en el diseño de sistemas de transcripción automatizada de audiencias, ya que permiten distribuir las cargas de procesamiento de manera eficiente entre los nodos de entrada (clientes) y los nodos de procesamiento central (servidores). En el contexto de sistemas ASR judiciales, el cliente captura el audio en tiempo real, lo envía mediante protocolos de transmisión de baja latencia (como WebRTC o WebSockets) y el servidor procesa los datos con motores de reconocimiento de voz basados en inteligencia artificial. Esta arquitectura asegura que el procesamiento intensivo de cómputo se realice en el backend, reduciendo la carga de los dispositivos de captura y garantizando escalabilidad (Kasakowskij, 2025).

La arquitectura distribuida se convierte en un complemento natural del modelo cliente-servidor al permitir que múltiples servidores trabajen en paralelo en tareas de preprocesamiento de audio, segmentación de hablantes y transcripción. Vera (2021) utilizó un esquema basado en FastAPI como backend y RabbitMQ/Celery como gestores de colas de tareas, lo que permitió procesar audios de manera asíncrona y mejorar la capacidad de respuesta del sistema.

\begin{table}[H]
\centering
\caption{Componentes clave en modelos cliente-servidor y arquitecturas distribuidas para sistemas ASR judiciales}
\begin{tabular}{|p{2.5cm}|p{3.5cm}|p{3cm}|p{3cm}|}
\hline
\rowcolor{gray!20}
\textbf{Componente} & \textbf{Descripción} & \textbf{Herramientas} & \textbf{Beneficio Principal} \\
\hline
Cliente & Captura de audio y envío al servidor en tiempo real & WebRTC, WebSockets & Baja latencia y sincronización en vivo \\
\hline
Servidor de Aplicación & Procesa el audio recibido, gestiona colas y coordina tareas & FastAPI, Celery, Redis & Escalabilidad y procesamiento asíncrono \\
\hline
Motor de ASR & Realiza la transcripción automática utilizando modelos IA & Deepgram Nova-3, faster-whisper & Alta precisión en reconocimiento de voz \\
\hline
Balanceador de Carga & Distribuye solicitudes entre múltiples nodos de servidor & NGINX, Kubernetes Ingress & Reducción de cuellos de botella \\
\hline
Seguridad y Privacidad & Protege la información sensible transmitida y almacenada & TLS 1.3, JWT, OAuth 2.0 & Confidencialidad e integridad de datos \\
\hline
\end{tabular}
\label{tab:cliente_servidor}
\end{table}

\subsection{Modelos Multimodales y Fusión de Señales}

Los modelos multimodales han surgido como una alternativa avanzada para mejorar la precisión de los sistemas de transcripción en entornos judiciales. Estos modelos combinan diferentes tipos de datos —por ejemplo, audio, video y texto— para generar representaciones más robustas y completas. Chaloupka, Palecek y Cerva (2021) demostraron que la integración de información visual junto con las señales acústicas mejora la segmentación de turnos de habla y reduce errores en entornos de ruido.

El enfoque multimodal puede implementarse mediante diferentes técnicas de fusión: temprana (\textit{early fusion}), tardía (\textit{late fusion}) e híbrida. Baltrušaitis et al. (2019) han demostrado que la fusión temprana es más eficaz en tareas que requieren correlación directa entre modalidades, mientras que la fusión tardía es más robusta frente a la pérdida de datos en alguna de las fuentes de señal.

\begin{table}[H]
\centering
\caption{Tipos de fusión de señales en modelos multimodales para ASR judicial}
\begin{tabular}{|p{2.5cm}|p{3cm}|p{2.5cm}|p{2.5cm}|p{2.5cm}|}
\hline
\rowcolor{gray!20}
\textbf{Tipo de Fusión} & \textbf{Descripción Técnica} & \textbf{Ventajas} & \textbf{Limitaciones} & \textbf{Aplicación en Justicia} \\
\hline
Fusión Temprana & Combina señales en la etapa de entrada & Mayor correlación entre modalidades & Requiere datos completos de todas las modalidades & Transcripción en audiencias con video de alta calidad \\
\hline
Fusión Tardía & Procesa cada modalidad de forma independiente & Robusta frente a pérdida de datos & Puede perder relaciones temporales finas & Audiencias con video parcial o señal de baja calidad \\
\hline
Fusión Híbrida & Combina fusión temprana y tardía & Balance entre robustez y correlación temporal & Complejidad computacional elevada & Sistemas judiciales de alta exigencia \\
\hline
Fusión Contextual & Integra metadatos y contexto legal & Enriquece el resultado con información semántica & Requiere bases de datos bien estructuradas & Generación automática de actas con identificación de intervinientes \\
\hline
\end{tabular}
\label{tab:fusion}
\end{table}

\section{Métricas y Parámetros de Evaluación}

La evaluación de sistemas de reconocimiento automático del habla en contextos judiciales exige el uso de métricas robustas que permitan medir no solo la precisión de las transcripciones, sino también su utilidad práctica en la producción de actas oficiales. Entre las métricas más empleadas se encuentra el Word Error Rate (WER), que cuantifica los errores de sustitución, inserción y omisión en relación con el número total de palabras de referencia:

\begin{equation}
\text{WER} = \frac{S + D + I}{N} \times 100
\end{equation}

donde $S$ representa sustituciones, $D$ eliminaciones, $I$ inserciones y $N$ el número total de palabras de referencia (Garneau y Bolduc, 2024).

\begin{table}[H]
\centering
\caption{Métricas y parámetros de evaluación en sistemas ASR judiciales}
\begin{tabular}{|p{3cm}|p{3.5cm}|p{3cm}|p{3cm}|}
\hline
\rowcolor{gray!20}
\textbf{Métrica / Parámetro} & \textbf{Definición} & \textbf{Método} & \textbf{Aplicación Judicial} \\
\hline
Word Error Rate (WER) & Porcentaje de errores de palabras respecto al total de referencia & $(S+D+I)/N\times100$ & Medición de precisión general de la transcripción \\
\hline
Character Error Rate (CER) & Porcentaje de errores a nivel de caracteres & Similar a WER pero a nivel de caracteres & Útil para lenguas con alta complejidad morfológica \\
\hline
BLEU Score & Evalúa correspondencia entre texto generado y referencia & Cálculo basado en n-gramas coincidentes & Útil en traducción automática de audiencias multilingües \\
\hline
MOS (Mean Opinion Score) & Medida subjetiva de calidad percibida & Escala de 1 a 5 evaluada por expertos & Validación de inteligibilidad y naturalidad en actas \\
\hline
Latencia (ms) & Tiempo promedio de procesamiento por unidad de audio & Medición en pruebas de estrés & Determina viabilidad para audiencias en tiempo real \\
\hline
Uso de GPU/CPU & Consumo de recursos computacionales & Monitoreo en entornos de prueba & Requisito para estimar costos y escalabilidad \\
\hline
\end{tabular}
\label{tab:metricas}
\end{table}

\subsection{Word Error Rate (WER) y Sonnex Distance}

La Word Error Rate (WER) es la métrica de referencia más utilizada en la evaluación de sistemas ASR, ya que permite cuantificar de manera objetiva el nivel de coincidencia entre la transcripción automática y la transcripción de referencia (\textit{ground truth}). Su fórmula estándar es:

\begin{equation}
\text{WER} = \frac{S + D + I}{N} \times 100
\end{equation}

Un valor de WER del 0\,\% indica una transcripción perfecta, mientras que valores superiores al 30\,\% suelen considerarse inaceptables en contextos críticos como el judicial (Harrington, 2023). En investigaciones recientes se reportan WER promedios del 15\,\% para sistemas como Amazon Transcribe y del 18\,\% para Whisper en condiciones de audio limpio, pero estos valores se incrementan hasta 50\,\% en audios forenses con ruido de fondo (Loakes, 2024). En contraste, servicios comerciales optimizados para streaming como Deepgram Nova-3 reportan WER de aproximadamente 8,4\,\% en español en condiciones estándar, con la ventaja adicional de operar en tiempo real con latencias inferiores a 300\,ms.

A diferencia del WER, que se limita a medir coincidencias léxicas, la \textbf{Sonnex Distance} es una métrica de carácter fonético que compara la similitud de las secuencias de sonido entre el texto de referencia y la transcripción automática. Este enfoque permite identificar errores que no son necesariamente léxicos, como variaciones de género, número o errores en homófonos (Garneau y Bolduc, 2024).

\begin{table}[H]
\centering
\caption{Métricas de evaluación léxica y fonética en sistemas ASR}
\begin{tabular}{|p{3cm}|p{3cm}|p{4cm}|p{3cm}|}
\hline
\rowcolor{gray!20}
\textbf{Métrica} & \textbf{Aspecto Evaluado} & \textbf{Fórmula / Método} & \textbf{Aplicación Judicial} \\
\hline
WER & Coincidencia léxica de palabras & $(S+D+I)/N \times 100$ & Detecta errores de sustitución, inserción y omisión en actas de audiencia \\
\hline
Sonnex Distance & Similitud fonética & Distancia de edición fonética ponderada & Identifica errores de homófonos, género y número; útil en entornos multilingües \\
\hline
\end{tabular}
\label{tab:wer_sonnex}
\end{table}

\subsection{Métricas de Latencia y Tiempo de Procesamiento}

En los sistemas ASR, la latencia y el tiempo de procesamiento son métricas críticas que determinan la viabilidad de su uso en entornos de tiempo real. La latencia se define como el intervalo de tiempo que transcurre entre la entrada de la señal de audio y la generación de la transcripción correspondiente (Macháček et al., 2023). Su cálculo es:

\begin{equation}
\text{Latencia (ms)} = t_{\text{salida}} - t_{\text{entrada}}
\end{equation}

donde $t_{\text{entrada}}$ corresponde al tiempo de captura del primer frame de audio y $t_{\text{salida}}$ al instante en que se produce el primer token de texto. El tiempo de procesamiento puede expresarse mediante el Real-Time Factor (RTF):

\begin{equation}
\text{RTF} = \frac{t_{\text{procesamiento}}}{t_{\text{audio}}}
\end{equation}

Un valor de RTF $< 1$ indica que el sistema es más rápido que el tiempo real, lo cual es deseable para aplicaciones de audiencias virtuales.

\begin{table}[H]
\centering
\caption{Métricas de evaluación de desempeño en sistemas ASR}
\begin{tabular}{|p{2.5cm}|p{3.5cm}|p{3cm}|p{3cm}|}
\hline
\rowcolor{gray!20}
\textbf{Métrica} & \textbf{Definición} & \textbf{Parámetro Clave} & \textbf{Valor Óptimo para Audiencias} \\
\hline
Latencia & Tiempo entre la captura del audio y la salida de la transcripción & $t_{\text{salida}} - t_{\text{entrada}}$ & $< 300$ ms \\
\hline
Tiempo de Procesamiento (RTF) & Proporción entre tiempo de cómputo y duración del audio procesado & $t_{\text{procesamiento}} / t_{\text{audio}}$ & $< 1{,}0$ \\
\hline
Throughput & Volumen de datos procesados por unidad de tiempo & Palabras/s o frames/s & $\geq$ 150 wpm \\
\hline
\end{tabular}
\label{tab:latencia}
\end{table}

\subsection{Parámetros Acústicos y de Entrenamiento}

El desempeño de un sistema ASR depende en gran medida de la correcta selección de parámetros acústicos y de la estrategia de entrenamiento aplicada al modelo. Los parámetros acústicos más comunes son los coeficientes cepstrales en las frecuencias de Mel (MFCCs), los filtros log-Mel y los espectrogramas de potencia (Gulati et al., 2020). La extracción de parámetros MFCC se resume en la siguiente secuencia matemática:

\begin{equation}
c_n = \text{DCT}\left[\log\left(M \cdot |FFT(x(t) \cdot w(t))|^2\right)\right]
\end{equation}

donde $x(t)$ es la señal de entrada, $w(t)$ es la ventana de análisis (ventana Hamming de 25 ms), $FFT$ es la Transformada Rápida de Fourier, $M$ es el banco de filtros de Mel y $DCT$ es la Transformada Discreta del Coseno.

\begin{table}[H]
\centering
\caption{Parámetros acústicos y de entrenamiento en modelos ASR}
\begin{tabular}{|p{2.5cm}|p{3.5cm}|p{2.5cm}|p{3.5cm}|}
\hline
\rowcolor{gray!20}
\textbf{Categoría} & \textbf{Parámetro} & \textbf{Valor Típico} & \textbf{Importancia en el Modelo} \\
\hline
\multirow{3}{2.5cm}{Acústicos} & Tamaño de ventana (frame) & 25 ms (solapamiento 10 ms) & Determina la resolución temporal del análisis espectral \\
\cline{2-4}
& Número de filtros Mel & 40 filtros & Captura de energía en bandas de frecuencia perceptual \\
\cline{2-4}
& Coeficientes MFCC & 13-39 coeficientes & Representación compacta de la envolvente espectral \\
\hline
\multirow{4}{2.5cm}{Entrenamiento} & Dataset de entrenamiento & $>$10\,000 h de audio & Mejora la capacidad de generalización \\
\cline{2-4}
& Tasa de aprendizaje & 1e-4 a 1e-3 & Controla la convergencia de la red neuronal \\
\cline{2-4}
& Tamaño del batch & 16-64 ejemplos por iteración & Impacta en estabilidad y tiempo de entrenamiento \\
\cline{2-4}
& Épocas de entrenamiento & 50-200 & Define la cantidad de iteraciones sobre el dataset \\
\hline
\end{tabular}
\label{tab:parametros}
\end{table}

\subsection{Validación Estadística}

La validación estadística constituye un componente esencial en el desarrollo de sistemas ASR, ya que garantiza que los resultados obtenidos no sean producto del azar. Un enfoque común es el uso de pruebas t para muestras relacionadas, cuando se cuenta con pares de resultados antes y después de implementar el sistema:

\begin{equation}
t = \frac{\bar{d}}{s_d / \sqrt{n}}
\end{equation}

donde $\bar{d}$ es la media de las diferencias entre pares, $s_d$ es la desviación estándar de las diferencias y $n$ es el número de pares de observaciones. Si el valor de $p < 0{,}05$, se rechaza la hipótesis nula.

En investigaciones con varias condiciones experimentales, se utiliza el Análisis de Varianza (ANOVA):

\begin{equation}
F = \frac{MS_{\text{entre}}}{MS_{\text{dentro}}}
\end{equation}

donde $MS_{\text{entre}}$ representa la variación media entre grupos y $MS_{\text{dentro}}$ la variación media dentro de cada grupo. En contextos donde no se cumplen supuestos de normalidad, se emplean métodos no paramétricos como la prueba de Wilcoxon para muestras emparejadas o la prueba de Kruskal-Wallis para comparaciones de múltiples grupos (Larrea, 2024).

\section{Desarrollo del Proyecto}

El desarrollo del presente proyecto se orienta a la implementación de un sistema automatizado de transcripción de audiencias judiciales basado en inteligencia artificial, cuyo propósito es optimizar la calidad y la celeridad de la producción de actas judiciales en el Distrito Judicial de Cusco. El sistema se estructura en dos etapas de implementación: la primera etapa entrega un producto mínimo viable (MVP) con transcripción en tiempo real, detección de múltiples hablantes, edición asistida y generación de actas; la segunda etapa incorpora mejoras de precisión, procesamiento por lotes, flujo de aprobación de actas y exportación documental en formatos DOCX y PDF.

El sistema integra únicamente dos servicios externos de inteligencia artificial: \textbf{Deepgram Nova-3} para el reconocimiento automático del habla en tiempo real y \textbf{Claude Sonnet 4} (Anthropic) para el mejoramiento lingüístico del texto y la generación de actas con formato oficial. Los demás componentes ---base de datos, procesamiento de audio, diccionario jurídico, editor de texto, exportación documental--- operan localmente en la infraestructura del sistema.

\subsection{Arquitectura de Software}

La arquitectura de software ha sido diseñada bajo un modelo cliente-servidor distribuido, asegurando la interoperabilidad de los componentes y la posibilidad de integrarse con los sistemas de gestión documental existentes en el Poder Judicial. El sistema se compone de seis servicios principales orquestados mediante Docker Compose: (1) frontend Next.js 14; (2) backend FastAPI; (3) PostgreSQL 16; (4) Redis; (5) Celery worker; y (6) nginx como proxy inverso.

\subsection{Stack Tecnológico}

\begin{table}[H]
\centering
\small
\caption{Stack tecnológico del sistema}
\begin{tabular}{|p{2.2cm}|p{3.2cm}|p{7cm}|}
\hline
\rowcolor{gray!20}
\textbf{Capa} & \textbf{Tecnología} & \textbf{Justificación} \\
\hline
\multirow{5}{2.2cm}{Frontend} & Next.js 14 & Framework React con App Router y renderizado del lado del servidor \\
\cline{2-3}
& TypeScript & Tipado estático para reducir errores en desarrollo \\
\cline{2-3}
& TipTap (ProseMirror) & Editor extensible para el canvas de transcripción \\
\cline{2-3}
& Zustand & Gestor de estado ligero y reactivo \\
\cline{2-3}
& wavesurfer.js & Reproductor de audio con forma de onda sincronizado \\
\hline
\multirow{4}{2.2cm}{Backend} & FastAPI & Framework web asíncrono con documentación OpenAPI automática \\
\cline{2-3}
& SQLAlchemy async + asyncpg & ORM asíncrono con driver nativo para PostgreSQL \\
\cline{2-3}
& Alembic & Gestión de migraciones de base de datos \\
\cline{2-3}
& Celery + Redis & Cola de tareas para procesamiento por lotes \\
\hline
Base de datos & PostgreSQL 16 & Base de datos relacional con soporte JSONB \\
\hline
\multirow{2}{2.2cm}{APIs Externas} & Deepgram Nova-3 & ASR en tiempo real con diarización nativa \\
\cline{2-3}
& Claude Sonnet 4 & Mejoramiento lingüístico y generación de actas \\
\hline
\multirow{2}{2.2cm}{Proc. local} & faster-whisper large-v3 & ASR por lotes con CTranslate2 \\
\cline{2-3}
& Pyannote Audio 3.1 & Diarización de hablantes por lotes \\
\hline
Infraestructura & Docker Compose + nginx & Orquestación de servicios y proxy inverso con TLS \\
\hline
Exportación & python-docx / weasyprint & Generación de DOCX y PDF con formato oficial \\
\hline
\end{tabular}
\label{tab:stack}
\end{table}

\subsection{Servicios Externos de Inteligencia Artificial}

El sistema utiliza exclusivamente dos servicios externos de inteligencia artificial:

\begin{enumerate}
\item \textbf{Deepgram Nova-3:} Servicio ASR basado en la nube que ofrece transcripción en tiempo real mediante conexión WebSocket. Se configura con idioma español latinoamericano (\texttt{es-419}), diarización activada (\texttt{diarize=true}), formato inteligente (\texttt{smart\_format=true}) y hasta 100 keyterms jurídicos por sesión. Cada resultado incluye identificador de hablante, nivel de confianza por palabra e indicador de resultado final.

\item \textbf{Claude Sonnet 4 (Anthropic):} Modelo de lenguaje utilizado para: (a) el mejoramiento del texto transcrito ---corrección de puntuación, capitalización de términos jurídicos (Juez, Fiscal, Señoría), signos de interrogación, sin añadir palabras no pronunciadas---; y (b) la generación del acta judicial formal a partir del contenido del canvas y los metadatos de la audiencia, con prompts específicos para el formato oficial del Poder Judicial del Perú.
\end{enumerate}

Todos los demás componentes ---diccionario jurídico, procesamiento de audio, editor de texto, base de datos, exportación documental, procesamiento por lotes--- operan localmente sin depender de servicios externos adicionales.

\subsection{Interfaz de Usuario (Frontend -- Next.js 14)}

El frontend es una aplicación web desarrollada con Next.js 14 y React 18 en TypeScript, que comprende los siguientes componentes principales:

\begin{itemize}
\item \textbf{Canvas de transcripción:} Editor TipTap con extensiones personalizadas: \texttt{SpeakerNode} para etiquetas de hablante, \texttt{SegmentMark} para segmentos con timestamps, \texttt{LowConfidenceMark} para palabras de baja confianza (fondo amarillo si confianza $< 0{,}7$), \texttt{ProvisionalNode} para texto provisional (gris) y \texttt{BookmarkNode} para marcadores.
\item \textbf{Panel de hablantes:} Lista de hablantes detectados con selector de rol (Juez, Fiscal, Defensa, Imputado, Agraviado, Testigo), etiqueta y color diferenciado.
\item \textbf{Reproductor de audio:} Implementado con wavesurfer.js: forma de onda, controles de reproducción, velocidad y salto por timestamp. Un clic en un segmento del canvas mueve el reproductor al timestamp correspondiente.
\item \textbf{Panel de marcadores:} Lista cronológica de marcadores con nota y timestamp; la selección navega al punto correspondiente en el canvas y en el reproductor.
\item \textbf{Sistema de sugerencias:} Popovers para correcciones del diccionario jurídico y análisis contextual de palabras de baja confianza. El digitador acepta con Tab o rechaza con Esc.
\item \textbf{Editor de acta:} Editor TipTap dedicado para la edición del acta generada, con estructura según el formato oficial del Poder Judicial.
\item \textbf{Barra de estado:} Conteo de palabras, tiempo transcurrido, estado de conexión Deepgram y número de segmentos.
\end{itemize}

\subsection{Captura y Almacenamiento de Audio}

La captura de audio se realiza en el navegador mediante la Web Audio API a 16\,kHz mono, con chunks cada 250\,ms enviados por WebSocket al backend. En el servidor, un módulo de grabación abre un archivo WAV (16\,bit, 16\,kHz, mono) por audiencia. Al finalizar la sesión, el archivo WAV se cierra y se actualiza \texttt{audio\_path} y duración en la base de datos, garantizando la preservación íntegra del registro sonoro para consulta posterior y procesamiento por lotes.

Para la segunda etapa, el backend acepta archivos pregrabados en formatos WAV, MP3, M4A y OGG, normalizando la frecuencia de muestreo mediante FFmpeg. Estudios previos han demostrado que el preprocesamiento puede reducir el WER hasta en un 20\,\% en entornos con ruido (Harrington, 2023).

\subsection{Comunicación en Tiempo Real (WebSocket)}

El protocolo WebSocket establece un canal bidireccional persistente entre el frontend y el backend. El endpoint \texttt{/ws/transcripcion/\{audiencia\_id\}} requiere autenticación JWT y valida el permiso del usuario. Se configura con timeout de 3600 segundos en nginx para audiencias de hasta 5 horas. Los mensajes incluyen: del cliente al servidor, chunks de audio (base64 PCM); del servidor al cliente, segmentos transcritos (\texttt{transcript}) y sugerencias del diccionario (\texttt{suggestion}).

\subsection{Servidor Backend (FastAPI)}

El servidor backend, desarrollado con FastAPI sobre Uvicorn (ASGI), se estructura en paquetes: \texttt{api} (endpoints REST), \texttt{ws} (handlers WebSocket), \texttt{services} (lógica de negocio), \texttt{models} (entidades SQLAlchemy) y \texttt{schemas} (validación Pydantic). Los servicios principales incluyen: autenticación JWT, API REST para audiencias, segmentos, hablantes, marcadores, actas y usuarios; \texttt{DeepgramStreamingService} para la conexión con Deepgram; \texttt{RealTimeEnhancementService} para invocación a Claude; \texttt{LegalDictionary} para el diccionario jurídico; y tareas Celery para procesamiento por lotes.

\subsection{Módulo de Transcripción en Tiempo Real}

El flujo de transcripción sigue la secuencia: (1) el digitador inicia la transcripción; (2) el frontend envía chunks de audio cada 250\,ms por WebSocket; (3) el backend reenvía a Deepgram Nova-3 y recibe resultados con speaker, confianza y palabras; (4) el servicio de consolidación agrupa segmentos del mismo hablante hasta formar frases completas; (5) el segmento consolidado se envía a Claude para mejoramiento y al diccionario jurídico; (6) el segmento final se persiste en PostgreSQL y se envía al canvas; (7) las sugerencias del diccionario se envían como mensajes \texttt{suggestion}.

\subsection{Módulo de Mejoramiento de Texto con IA}

El servicio \texttt{RealTimeEnhancementService} invoca a Claude Sonnet 4 para cada segmento, con las reglas: mejorar puntuación y capitalizar términos jurídicos; reemplazar palabras mal transcritas (sustitución 1:1); \textbf{nunca añadir palabras} no pronunciadas ni completar frases; proporcionar contexto de los últimos 25 segmentos. La regla de no añadir palabras es un principio fundamental: toda sugerencia es exclusivamente un reemplazo, nunca una inserción.

\subsection{Módulo de Diccionario Jurídico}

El diccionario jurídico opera localmente sin dependencia de servicios externos. Su base de datos (\texttt{legal\_terms.json}) contiene términos jurídicos con variantes erróneas, categoría y contexto. Emplea Soundex adaptado al español y distancia de Levenshtein para búsqueda fuzzy (por ejemplo, ``sobresemiento'' $\to$ ``sobreseimiento''). Se configuran hasta 100 keyterms jurídicos por sesión en Deepgram, organizados por categorías: procesal penal, delitos, partes procesales, normativa y terminología regional de Cusco. Objetivo de rendimiento: $< 50$\,ms por segmento.

\subsection{Módulo de Gestión de Hablantes}

La diarización se implementa en dos niveles: en tiempo real, Deepgram Nova-3 asigna identificadores (SPEAKER\_00, SPEAKER\_01, etc.) y el panel de hablantes permite asociar cada uno con un rol procesal (Juez, Fiscal, Defensa, Imputado, Agraviado, Testigo, Perito, Secretario); en procesamiento por lotes (segunda etapa), Pyannote Audio 3.1 realiza diarización más precisa sobre el archivo WAV completo, alineando temporalmente con la transcripción de faster-whisper.

\subsection{Módulo de Procesamiento por Lotes}

Implementado en la segunda etapa, permite reprocesar el audio con mayor precisión: \textbf{faster-whisper large-v3} transcribe el WAV completo con \texttt{beam\_size=5}, \texttt{word\_timestamps=True} y prompt inicial con terminología judicial; \textbf{Pyannote Audio 3.1} identifica turnos de habla (requiere GPU NVIDIA RTX 3060 12\,GB mínimo). El servicio de alineación cruza palabras de faster-whisper con segmentos de Pyannote. Los resultados se comparan con los segmentos del streaming; cuando la diferencia supera un umbral configurable, se genera una propuesta que el digitador acepta o rechaza en vista de comparación. Los segmentos con \texttt{editado\_por\_usuario=true} no reciben propuestas. La tarea se ejecuta mediante Celery con Redis.

\subsection{Módulo de Generación de Actas}

En la segunda etapa, el digitador genera el acta pulsando un botón. El sistema captura un snapshot del canvas y metadatos de la audiencia (expediente, juzgado, tipo, fecha, hablantes con roles). El contenido se envía a Claude Sonnet 4 con prompt según formato oficial (Formato A para juzgados unipersonales, Formato B para salas de apelaciones). El acta se almacena con versión y estado ``borrador'', y se carga en el editor dedicado. El flujo de aprobación contempla: borrador $\to$ en\_revisión $\to$ aprobada $\to$ exportada. Solo el rol supervisor puede aprobar, registrando quién aprobó y la fecha.

\subsection{Módulo de Exportación Documental}

La exportación, implementada en la segunda etapa, genera: \textbf{DOCX} mediante python-docx con plantilla de estilos oficiales (encabezado, títulos, cuerpo, firmas) y datos insertados programáticamente; \textbf{PDF} mediante weasyprint (HTML a PDF). La exportación está disponible cuando el acta alcanza estado ``aprobada''. Cada exportación se registra en log de auditoría.

\subsection{Modelo de Datos}

\begin{table}[H]
\centering
\small
\caption{Entidades principales del modelo de datos}
\begin{tabular}{|p{2.2cm}|p{4cm}|p{6.5cm}|}
\hline
\rowcolor{gray!20}
\textbf{Entidad} & \textbf{Campos clave} & \textbf{Descripción} \\
\hline
Audiencia & expediente, juzgado, tipo\_audiencia, instancia, fecha, estado, audio\_path & Sesión judicial con metadatos completos \\
\hline
Segmento & audiencia\_id, speaker\_id, texto\_ia, texto\_mejorado, texto\_editado, timestamps, confianza, orden, editado\_por\_usuario, palabras\_json (JSONB) & Unidad de transcripción con trazabilidad \\
\hline
Hablante & audiencia\_id, speaker\_id, rol, etiqueta\_canvas, color & Rol procesal de cada hablante detectado \\
\hline
Marcador & audiencia\_id, nota, timestamp\_audio & Punto de interés con nota del digitador \\
\hline
Acta & audiencia\_id, contenido\_llm, contenido\_final, version, estado, aprobada\_por & Acta con historial y flujo de aprobación \\
\hline
Usuario & email, nombre, rol, password\_hash, activo & Roles: admin, transcriptor, supervisor \\
\hline
\end{tabular}
\label{tab:modelo_datos}
\end{table}

\subsection{Reglas de Negocio del Sistema}

\begin{enumerate}
\item \textbf{Regla de no sobreescritura:} Si \texttt{editado\_por\_usuario=true}, ningún proceso automático puede modificar ese segmento. El digitador siempre tiene el control final.
\item \textbf{Regla de corrección 1:1:} Toda sugerencia ---diccionario, análisis contextual o Claude--- es exclusivamente un reemplazo de una palabra por otra; no se insertan palabras nuevas.
\item \textbf{Umbrales de confianza:} Sugerencias de reemplazo solo para palabras con confianza $< 0{,}85$. Las palabras correctas del diccionario se excluyen.
\item \textbf{Consolidación de frases:} Segmentos de menos de 5 palabras no se procesan como finales, salvo respuestas cortas válidas (sí, no, niego). Buffers con más de 20 palabras que terminan en conectores se consultan a Claude.
\end{enumerate}

\subsection{Seguridad del Sistema}

\begin{itemize}
\item \textbf{Autenticación:} JWT con access token y refresh token; validación de permiso por audiencia en el WebSocket.
\item \textbf{Contraseñas:} Hashing con bcrypt (12 rondas).
\item \textbf{Privacidad:} Solo texto se envía a Claude, nunca audio.
\item \textbf{Transporte:} HTTPS en producción con nginx como terminación TLS.
\item \textbf{Configuración:} Variables sensibles en archivo \texttt{.env} no versionado.
\end{itemize}

\subsection{Requerimientos No Funcionales}

\begin{table}[H]
\centering
\small
\caption{Métricas de rendimiento del sistema}
\begin{tabular}{|p{5cm}|p{3cm}|p{5cm}|}
\hline
\rowcolor{gray!20}
\textbf{Métrica} & \textbf{Objetivo} & \textbf{Justificación} \\
\hline
Latencia de primera transcripción & $< 2$ s & Feedback en tiempo real para el digitador \\
\hline
Diccionario jurídico & $< 50$ ms & No bloquear el flujo de segmentos \\
\hline
Conexión WebSocket & Timeout 3600 s & Audiencias de 2--5 horas sin corte \\
\hline
Generación de acta & $< 60$ s & Experiencia fluida al finalizar \\
\hline
Exportación DOCX/PDF & $< 10$ s & Descarga sin espera excesiva \\
\hline
Disponibilidad & 99\% mensual & En horario de audiencias (8:00--18:00) \\
\hline
\end{tabular}
\label{tab:rendimiento}
\end{table}

Compatibilidad: Chrome o Edge para captura de audio; servidor Linux con Docker; Python 3.11, Node 18+. Escalabilidad: múltiples audiencias concurrentes. Mantenibilidad: código modular, migraciones con Alembic, documentación OpenAPI.

\subsection{Alcance por Etapa de Implementación}

\begin{table}[H]
\centering
\small
\caption{Alcance funcional por etapa de implementación}
\begin{tabular}{|c|p{2.5cm}|p{9cm}|}
\hline
\rowcolor{gray!20}
\textbf{Etapa} & \textbf{Alcance} & \textbf{Funcionalidades principales} \\
\hline
1 (MVP) & Producto mínimo viable & Infraestructura y base de datos; autenticación y roles; transcripción en tiempo real con Deepgram Nova-3 y diarización; captura y almacenamiento de audio WAV; consolidación por hablante y mejoramiento con Claude; diccionario jurídico con fuzzy matching; canvas TipTap con hablantes, marcadores y sugerencias; frases estándar; generación de acta borrador; despliegue con Docker Compose \\
\hline
2 & Precisión, batch, acta y exportación & Mejoras de precisión: regla ``no añadir palabras'', corrección 1:1, umbrales de confianza; procesamiento por lotes (faster-whisper + Pyannote) con propuestas de mejora; editor de acta con versionado; flujo de aprobación por supervisor; exportación DOCX y PDF con formato oficial; estabilización y pruebas E2E \\
\hline
\end{tabular}
\label{tab:etapas}
\end{table}

\subsection{Flujo General del Sistema}

El flujo de datos completo es: (1) creación de la audiencia con metadatos; (2) inicio de captura de audio en el navegador; (3) transmisión por WebSocket al backend; (4) reenvío a Deepgram Nova-3 con diarización; (5) consolidación de segmentos por hablante; (6) mejoramiento con Claude Sonnet 4; (7) verificación contra el diccionario jurídico; (8) persistencia en PostgreSQL; (9) envío del segmento y sugerencias al canvas; (10) edición por el digitador; (11) generación del acta con Claude; (12) aprobación por supervisor; (13) exportación a DOCX/PDF.

\subsection{Metodología de Desarrollo de Software}

El desarrollo sigue un enfoque ágil basado en Scrum, con iteraciones cortas que concluyen con la entrega de un incremento potencialmente desplegable. Las funcionalidades se priorizan según la técnica MoSCoW (Must, Should, Could, Won't) y se validan mediante demos periódicas con los usuarios del Distrito Judicial.

\subsection{Componentes de la Sala de Audiencia}

La infraestructura física es un componente crítico para garantizar la calidad del audio y video. La sala está equipada con micrófonos SHURE MX418 D/C estratégicamente ubicados, cámaras Logitech Brio Ultra HD, y un sistema de sonido JBL, lo que asegura una captura clara de las intervenciones (Larrea, 2024). Además, se cuenta con un UPS para evitar interrupciones por fallas de energía y con mobiliario ergonómico que facilita la disposición ordenada de los equipos.

\section{Marco Conceptual}

El marco conceptual constituye la base semántica y teórica de la presente investigación, ya que precisa y define los términos técnicos, jurídicos y tecnológicos utilizados a lo largo del trabajo.

\subsection{Inteligencia Artificial (IA)}

La inteligencia artificial (IA) es un campo de la informática que busca desarrollar sistemas capaces de realizar tareas que normalmente requieren inteligencia humana, tales como el aprendizaje, el razonamiento, la planificación y el reconocimiento de patrones. La IA se divide en ramas como el aprendizaje automático (\textit{machine learning}) y el aprendizaje profundo (\textit{deep learning}), que permiten a los modelos mejorar su rendimiento a medida que se exponen a nuevos datos (Russell y Norvig, 2024).

En el contexto de esta investigación, la IA es el núcleo de la solución tecnológica, ya que habilita el reconocimiento de voz mediante Deepgram Nova-3, el mejoramiento lingüístico mediante Claude Sonnet 4 y el procesamiento por lotes mediante faster-whisper. Investigaciones recientes muestran que los modelos basados en \textit{deep learning} pueden alcanzar niveles de precisión superiores al 90\,\% en entornos controlados (Radford et al., 2022).

\subsection{Reconocimiento Automático del Habla}

El reconocimiento automático del habla (ASR) es una tecnología que convierte señales acústicas en texto mediante la combinación de modelos acústicos, modelos de lenguaje y algoritmos de decodificación. Este proceso suele involucrar la extracción de características como los coeficientes cepstrales en la frecuencia de Mel (MFCC), que sirven como insumo para redes neuronales profundas o arquitecturas híbridas DNN-HMM (Chaloupka et al., 2021). Para la presente investigación, el ASR constituye la piedra angular del sistema propuesto, pues sustituye la transcripción manual de audiencias, reduciendo tiempos y carga administrativa.

\subsection{Transcripción Judicial}

La transcripción judicial es el procedimiento que convierte en texto escrito las declaraciones orales, intervenciones y decisiones emitidas en una audiencia (Poder Judicial del Perú, 2023). Esta actividad es fundamental para garantizar el derecho al debido proceso, ya que el acta de audiencia se constituye en un documento oficial que respalda las actuaciones judiciales. En el Distrito Judicial de Cusco, la transcripción se realiza manualmente, lo que incrementa el riesgo de errores humanos y demora la disponibilidad de actas.

\subsection{Celeridad Procesal}

La celeridad procesal es un principio rector de la justicia que busca que los procesos se desarrollen en un plazo razonable y sin dilaciones indebidas. Este principio se encuentra consagrado en la Constitución Política del Perú y en tratados internacionales, como la Convención Americana sobre Derechos Humanos. La implementación de sistemas de transcripción automática contribuye directamente a este principio al reducir el tiempo de entrega de actas y resoluciones. Prescott (2020) sostiene que el uso de tecnologías de información en los tribunales es un catalizador para fortalecer la confianza ciudadana, siempre que se garantice la imparcialidad y se mantenga la integridad de los registros procesales.

\subsection{Eficiencia Institucional}

La eficiencia institucional se define como la capacidad de una organización para alcanzar sus objetivos con el mínimo de recursos, maximizando su productividad y efectividad. En el Poder Judicial, esta eficiencia se relaciona con la reducción de la sobrecarga administrativa y la agilización de procesos críticos como la transcripción de audiencias. La automatización de las transcripciones libera a los especialistas de tareas repetitivas, lo que incrementa la productividad institucional. En Europa, la digitalización de servicios públicos ha sido señalada como una vía para mejorar la agilidad y eficiencia del sector público, reduciendo costos y tiempos operativos (Bauer et al., 2025).

\subsection{FastAPI}

FastAPI es un framework web moderno de Python, diseñado para construir APIs de alto rendimiento con tipado estático y generación automática de documentación OpenAPI. Se ejecuta sobre servidores ASGI como Uvicorn, lo que le permite gestionar operaciones de entrada/salida de forma asíncrona y atender miles de conexiones concurrentes sin sacrificar rendimiento. En el sistema propuesto, FastAPI se emplea como el servidor backend que coordina la autenticación JWT, la gestión de audiencias, la persistencia en PostgreSQL mediante SQLAlchemy asíncrono, la ejecución de tareas diferidas con Celery y Redis, y la comunicación en tiempo real con el frontend, garantizando estabilidad y escalabilidad para un uso continuo en audiencias judiciales.

\subsection{WebSocket}

WebSocket es un protocolo que establece un canal de comunicación bidireccional persistente entre cliente y servidor, reduciendo la latencia en el intercambio de datos (Fette y Melnikov, 2011). Esta característica es crucial en entornos judiciales, donde la inmediatez en la transcripción de declaraciones puede incidir en la toma de decisiones. La literatura técnica respalda su uso en aplicaciones críticas de baja latencia, como telemedicina y monitoreo en tiempo real.

\subsection{Deepgram}

Deepgram es un servicio de reconocimiento automático del habla basado en la nube que ofrece transcripción en tiempo real y por lotes mediante modelos de inteligencia artificial propietarios. Su modelo Nova-3 proporciona alta precisión en múltiples idiomas, incluyendo español, con soporte nativo para detección de hablantes (\textit{diarization}), puntuación automática e indicadores de confianza por palabra. En el sistema propuesto, Deepgram Nova-3 se utiliza como motor de transcripción en tiempo real (\textit{streaming}), conectado directamente desde el frontend mediante una conexión WebSocket a través de su SDK oficial.

\subsection{Next.js}

Next.js es un framework de React para el desarrollo de aplicaciones web, que ofrece renderizado del lado del servidor, generación estática y enrutamiento basado en el sistema de archivos. Desarrollado por Vercel, Next.js permite construir interfaces de usuario modernas, responsivas y de alto rendimiento. En el sistema propuesto, Next.js 14 se utiliza como la plataforma del frontend, proporcionando una interfaz web completa que incluye el dashboard de audiencias, el canvas de transcripción, la gestión de hablantes y la edición de actas.

\subsection{Whisper y faster-whisper}

Whisper es un modelo de ASR desarrollado por OpenAI, entrenado con 680\,000 horas de datos de audio multilingües. Su diseño basado en transformadores lo hace robusto frente a ruido de fondo, acentos y terminología especializada (Radford et al., 2022). En esta investigación, Whisper se emplea en su variante faster-whisper, una reimplementación optimizada mediante CTranslate2 que ofrece mayor velocidad de inferencia con menor consumo de memoria. Este motor se utiliza específicamente para el procesamiento por lotes (\textit{batch}) de audios pregrabados, complementando al servicio Deepgram Nova-3 que se encarga de la transcripción en tiempo real.

\subsection{Word Error Rate (WER)}

El WER es la métrica estándar para evaluar la precisión de los sistemas ASR. Se calcula mediante la fórmula:

\begin{equation}
\text{WER} = \frac{S + D + I}{N} \times 100
\end{equation}

donde $S$ representa sustituciones, $D$ eliminaciones, $I$ inserciones y $N$ el total de palabras de referencia (Morris et al., 2004). Un valor inferior al 10\,\% es aceptable en aplicaciones de alta exigencia, como la transcripción judicial (Garneau y Bolduc, 2024). Esta métrica será utilizada en la investigación para medir objetivamente el impacto del sistema propuesto.

\subsection{Claude (Anthropic)}

Claude es un modelo de lenguaje de gran escala desarrollado por Anthropic, diseñado para tareas de comprensión, generación y análisis de texto. En su versión Sonnet 4, ofrece un equilibrio entre precisión, velocidad y costo, siendo adecuado para aplicaciones en tiempo real. En el sistema propuesto, Claude se utiliza para dos funciones: el mejoramiento lingüístico de los segmentos transcritos (puntuación, capitalización de términos jurídicos, corrección ortográfica sin añadir palabras) y la generación del acta judicial formal a partir del contenido transcrito y los metadatos de la audiencia. A diferencia de los modelos de reconocimiento de voz, Claude opera exclusivamente sobre texto, lo que garantiza que nunca se transmite audio a servicios externos de procesamiento de lenguaje natural.

\subsection{TipTap}

TipTap es un framework de edición de texto enriquecido para aplicaciones web, construido sobre ProseMirror. Permite crear editores altamente personalizables mediante un sistema de extensiones modulares, donde cada funcionalidad (nodos, marcas, plugins) se añade de forma independiente. En el sistema propuesto, TipTap se utiliza para implementar el canvas de transcripción con extensiones personalizadas: nodos de hablante con etiquetas y colores, marcas de segmento con timestamps, indicadores de baja confianza y marcadores. Su arquitectura extensible permite integrar la recepción de segmentos en tiempo real por WebSocket directamente en el editor, ofreciendo una experiencia de edición fluida y profesional.

\subsection{PostgreSQL}

PostgreSQL es un sistema de gestión de bases de datos relacional de código abierto, reconocido por su robustez, extensibilidad y cumplimiento de estándares SQL. En su versión 16, incorpora soporte nativo para datos semiestructurados mediante el tipo JSONB, lo que permite almacenar información flexible como las palabras individuales con timestamps y confianza de cada segmento de transcripción. En el sistema propuesto, PostgreSQL almacena todas las entidades del modelo de datos: audiencias, segmentos, hablantes, marcadores, actas y usuarios.

\subsection{Docker}

Docker es una plataforma de virtualización a nivel de contenedores que permite empaquetar aplicaciones con todas sus dependencias en unidades aisladas y reproducibles. Docker Compose, su herramienta de orquestación, permite definir y ejecutar aplicaciones multicontenedor. En el sistema propuesto, Docker Compose orquesta seis servicios: frontend (Next.js), backend (FastAPI), PostgreSQL, Redis, Celery worker y nginx, asegurando un despliegue reproducible y consistente en cualquier entorno servidor.

\subsection{Redis}

Redis es un almacén de datos en memoria de código abierto, utilizado como base de datos, caché y broker de mensajes. En el sistema propuesto, Redis actúa como broker de mensajes para Celery, gestionando las colas de tareas del procesamiento por lotes (transcripción batch con faster-whisper y diarización con Pyannote). Su naturaleza en memoria permite una comunicación de baja latencia entre el backend y los workers de procesamiento.

% ============================================================================
% CAPÍTULO III
% ============================================================================
\chapter{METODOLOGÍA}

\section{Metodología de Investigación}

La metodología constituye el eje central de la investigación, pues permite organizar y ejecutar de manera sistemática las fases del proyecto, garantizando la validez y confiabilidad de los resultados obtenidos. En el caso del presente trabajo, se busca diseñar, implementar y evaluar un sistema automatizado de transcripción de audiencias judiciales empleando inteligencia artificial, específicamente mediante la integración de Deepgram Nova-3 para el reconocimiento automático del habla en tiempo real y Claude Sonnet 4 (Anthropic) para el mejoramiento lingüístico y la generación de actas judiciales. La metodología se fundamenta en el enfoque de investigación aplicada de desarrollo tecnológico, centrada en la construcción, prueba y validación de un prototipo funcional orientado a mejorar la precisión y eficiencia de la transcripción judicial en el Distrito Judicial de Cusco.

El proceso metodológico se estructura en ocho apartados: tipo y alcance de investigación, método y diseño, enfoque, unidad experimental, corpus de datos, técnicas de recolección, métricas de evaluación y metodología de desarrollo de software. El desarrollo del sistema sigue un enfoque ágil, que permite la integración iterativa de componentes tecnológicos, la validación temprana de resultados y la mejora continua durante las fases de diseño y evaluación.

\section{Tipo y Alcance de Investigación}

\subsection{Tipo de Investigación}

La investigación es de tipo aplicada y de desarrollo tecnológico, dado que busca generar una solución informática innovadora para un problema real del sistema judicial: la transcripción manual de audiencias que genera demoras, errores y sobrecarga administrativa. La investigación aplicada se caracteriza por la utilización del conocimiento científico para el diseño de soluciones prácticas, con el fin de mejorar procesos y optimizar recursos en un entorno específico (Hernández-Sampieri et al., 2022).

Asimismo, el estudio adopta un alcance tecnológico-explicativo, en tanto evalúa el efecto de la incorporación de inteligencia artificial sobre la precisión y eficiencia del proceso de transcripción, a partir de métricas técnicas como Word Error Rate (WER), Character Error Rate (CER) y Mean Opinion Score (MOS).

\subsection{Método de Investigación}

El método adoptado es experimental de tipo pre-experimental, en el sentido técnico del término, ya que se comparan los resultados obtenidos antes y después de la implementación del sistema, utilizando el mismo conjunto de audios judiciales como base. Se evalúa la precisión de las transcripciones manuales frente a las generadas automáticamente por el sistema, midiendo mejoras en tiempo y exactitud. Este enfoque es apropiado cuando no es posible contar con grupos de control, pero se busca obtener evidencia empírica del impacto de la intervención (Kerlinger y Lee, 2020).

En paralelo, el desarrollo del software sigue el método ágil bajo el marco de trabajo Scrum, lo que permite gestionar iterativamente el ciclo de vida del sistema.

\subsection{Diseño de Investigación}

El diseño de investigación es experimental de un solo grupo, orientado a comparar la precisión y eficiencia del proceso de transcripción judicial antes y después de la intervención tecnológica. Las métricas evaluadas incluyen WER, CER, MOS y latencia promedio, las cuales permiten cuantificar el desempeño del sistema en condiciones reales. Este diseño es ampliamente utilizado en estudios de adopción tecnológica, ya que permite validar la efectividad de la intervención en condiciones reales (Creswell y Creswell, 2018).

\subsection{Enfoque de Investigación}

El enfoque es cuantitativo de validación tecnológica, dado que los resultados se expresan mediante indicadores numéricos de desempeño del sistema. Las métricas se calculan utilizando métodos objetivos estandarizados en el campo del procesamiento de lenguaje natural (Garneau y Bolduc, 2024; Loakes, 2024). Complementariamente, se incluye un análisis descriptivo que contextualiza los resultados en el entorno judicial, permitiendo formular recomendaciones para su adopción institucional.

\section{Unidad de Análisis}

La unidad experimental está constituida por un conjunto de audios de audiencias judiciales obtenidos del Distrito Judicial de Cusco, que sirven como base para la evaluación del sistema. El corpus de datos incluye grabaciones anonimizadas de diferentes especialidades con una duración mínima de 20 minutos y participación de múltiples interlocutores.

\section{Población y Muestra de Estudio}

La validación del sistema se realiza mediante indicadores de precisión, latencia y calidad perceptual: WER (Word Error Rate), CER (Character Error Rate), MOS (Mean Opinion Score) y latencia promedio. Estas métricas permiten una evaluación objetiva y reproducible del desempeño del sistema, conforme a los estándares internacionales en reconocimiento automático del habla.



% ============================================================================
% REFERENCIAS
% ============================================================================
\chapter*{REFERENCIAS}
\addcontentsline{toc}{chapter}{Referencias}
\setlength{\parindent}{0pt}
\setlength{\hangindent}{1.25cm}

Ao, J., Wang, R., Zhou, L., Wei, Z., y Zhang, S. (2022). SpeechT5: Unified-modal encoder-decoder pre-training for spoken language processing. \textit{IEEE/ACM Transactions on Audio, Speech, and Language Processing, 30}, 1--12. \url{https://doi.org/10.18653/v1/2022.acl-long.393}

\vspace{6pt}
Baltrušaitis, T., Ahuja, C., y Morency, L. P. (2019). Multimodal machine learning: A survey and taxonomy. \textit{IEEE Transactions on Pattern Analysis and Machine Intelligence, 41}(2), 423--443. \url{https://doi.org/10.1109/TPAMI.2018.2798607}

\vspace{6pt}
Baevski, A., Zhou, H., Mohamed, A., y Auli, M. (2020). wav2vec 2.0: A framework for self-supervised learning of speech representations. \textit{Advances in Neural Information Processing Systems (NeurIPS), 33}, 12449--12460. \url{https://doi.org/10.48550/arXiv.2006.11477}

\vspace{6pt}
Bauer, M., Draghi, M., y otros. (2025). \textit{Boosting efficiency and quality in EU public services} (Occasional Paper No. 04/2025). ECIPE. \url{https://www.econstor.eu/bitstream/10419/315166/1/1920054340.pdf}

\vspace{6pt}
Bild, E., Redman, A., Newman, E. J., Muir, B. R., Tait, D., y Schwarz, N. (2021). Sound and credibility in the virtual court: Low audio quality leads to less favorable evaluations of witnesses and lower weighting of evidence. \textit{Law and Human Behavior, 45}(5), 481--495. \url{https://doi.org/10.1037/lhb0000466}

\vspace{6pt}
Cabrera Fernández, R. (2024). Automatización de procesos judiciales mediante inteligencia artificial: Implicancias para la celeridad procesal y la transparencia institucional. \textit{Revista Peruana de Derecho y Tecnología, 8}(1), 45--61. \url{https://doi.org/10.20318/eunomia.2024.9006}

\vspace{6pt}
Chaloupka, J., Palecek, J., y Cerva, P. (2021). Audio-visual broadcast transcription system using artificial neural networks. \textit{IEEE International Workshop on Electronics, Control, Measurement, Signals and their Application to Mechatronics (ECMSM)}, 1--6. \url{https://doi.org/10.1109/ECMSM51310.2021.9468830}

\vspace{6pt}
Comisión Económica para América Latina y el Caribe (CEPAL). (2025). \textit{Latin American Artificial Intelligence Index (ILIA 2025)}. \url{https://www.cepal.org/en/pressreleases/latin-america-and-caribbean-accelerate-adoption-artificial-intelligence}

\vspace{6pt}
Creswell, J. W., y Creswell, J. D. (2018). \textit{Research design: Qualitative, quantitative, and mixed methods approaches} (5.\textordfeminine{} ed.). SAGE.

\vspace{6pt}
Didena, T. (2024). Legal transcription services and access to justice: A qualitative analysis of court record-keeping in the Southern Nations, Nationalities, and Peoples' Region of Ethiopia. \textit{African Journal of Law and Society, 16}(1), 45--62. \url{https://doi.org/10.1177/13657127231217774}

\vspace{6pt}
Fette, I., y Melnikov, A. (2011). \textit{The WebSocket protocol} (RFC 6455). Internet Engineering Task Force (IETF). \url{https://doi.org/10.17487/RFC6455}

\vspace{6pt}
Garneau, G., y Bolduc, P. (2024). The state of commercial automatic French legal speech recognition systems and their impact on court reporters. \textit{Journal of Legal Technology Studies, 12}(1), 1--22. \url{https://arxiv.org/abs/2408.11940}

\vspace{6pt}
Gulati, A., Qin, J., Chiu, C.-C., Parmar, N., Zhang, Y., Yu, J., y Wu, Y. (2020). Conformer: Convolution-augmented transformer for speech recognition. \textit{Proceedings of Interspeech 2020}, 5036--5040. \url{https://doi.org/10.21437/Interspeech.2020-3015}

\vspace{6pt}
Harrington, J. (2023). Incorporating automatic speech recognition methods into the transcription of police-suspect interviews: Factors affecting automatic performance. \textit{Journal of Forensic Phonetics, 5}(2), 77--94. \url{https://doi.org/10.3389/fcomm.2023.1165233}

\vspace{6pt}
Hernández-Sampieri, R., Fernández Collado, C., y Baptista Lucio, P. (2022). \textit{Metodología de la investigación} (7.\textordfeminine{} ed.). McGraw-Hill.

\vspace{6pt}
Kasakowskij, M. (2025). Distributed architectures for real-time speech recognition: Applications in justice systems. \textit{International Journal of Cloud Computing and AI Systems, 14}(1), 50--68. \url{https://doi.org/10.1007/s44217-025-00568-6}

\vspace{6pt}
Kerlinger, F. N., y Lee, H. B. (2020). \textit{Foundations of behavioral research} (5.\textordfeminine{} ed.). Cengage.

\vspace{6pt}
Koenecke, A., Nam, A., Lake, E., Nudell, J., Quartey, M., Mengesha, Z., y Goel, S. (2020). Racial disparities in automated speech recognition. \textit{Proceedings of the National Academy of Sciences, 117}(14), 7684--7689. \url{https://doi.org/10.1073/pnas.1915768117}

\vspace{6pt}
Kostka, G., Steinacker, L., y Meckel, M. (2023). Under big brother's watchful eye: Cross-country attitudes toward facial recognition technology. \textit{Government Information Quarterly, 40}, 101761. \url{https://doi.org/10.1016/j.giq.2022.101761}

\vspace{6pt}
Larrea, P. (2024). \textit{Uso de la plataforma Google Meet y transcriptor IA en el proceso de trámite de las audiencias del módulo penal dentro del Poder Judicial de Junín en la sub sede central Huancayo 2024} [Tesis de maestría, Universidad Nacional del Centro del Perú]. \url{https://repositorio.uncp.edu.pe/handle/20.500.12894/12608}

\vspace{6pt}
Latif, S., Qureshi, S. A., Maqsood, M., y Ali, H. (2023). A survey of transformers in speech processing: Architectures, applications and challenges. \textit{IEEE Access, 11}, 12953--12974. \url{https://doi.org/10.1109/ACCESS.2023.3241234}

\vspace{6pt}
Legchenko, A., y Bondarenko, I. (2025). \textit{Adopting domain-specific knowledge in ASR systems}. Proceedings of the MathAI 2025 Conference. Novosibirsk State University. \url{https://openreview.net/pdf?id=dzrEXinn4N}

\vspace{6pt}
Loakes, D. (2024). Automatic speech recognition and the transcription of indistinct forensic audio: How do the new generation of systems fare? \textit{Frontiers in Communication, 9}, 1--12. \url{https://www.frontiersin.org/journals/communication/articles/10.3389/fcomm.2024.1281407/full}

\vspace{6pt}
Lyra, A., Barbosa, C. E., Salazar, H., Argôlo, M., Lima, Y., Motta, R., y Moreira de Souza, J. (2024). Automatic transcription systems: A game changer for court hearings. En \textit{Proceedings of the 16th International Joint Conference on Knowledge Discovery, Knowledge Engineering and Knowledge Management} (pp. 163--174). SCITEPRESS. \url{https://doi.org/10.5220/0012891400003838}

\vspace{6pt}
Macháček, J., Švec, J., y Žmolík, M. (2023). Low-latency automatic speech recognition systems for real-time applications: A performance evaluation. \textit{Journal of Real-Time Image Processing, 20}(3), 421--435. \url{https://doi.org/10.3389/frsip.2022.999457}

\vspace{6pt}
Morris, A., Maier, V., y Green, P. (2004). From WER and RIL to MER and WIL: Improved evaluation measures for connected speech recognition. \textit{IEEE Signal Processing Letters, 12}(10), 1--4.

\vspace{6pt}
Nascimento, F., Silva, R., y Duarte, L. (2024). Automated transcription in parliamentary sessions: Efficiency gains and legal considerations. \textit{Parliamentary Technology Review, 22}(3), 88--105. \url{https://doi.org/10.1177/ptr.2024.88}

\vspace{6pt}
OCDE. (2023). \textit{Principios de la OCDE sobre la inteligencia artificial}. Organización para la Cooperación y el Desarrollo Económicos. \url{https://oecd.ai}

\vspace{6pt}
Poder Judicial del Perú. (2023). \textit{Reglamento de audiencias virtuales y uso de medios tecnológicos en procesos judiciales}. Consejo Ejecutivo del Poder Judicial.

\vspace{6pt}
Prescott, J. J. (2020). Improving access to justice in state courts with platform technology. \textit{Vanderbilt Law Review, 73}(6), 1875--1928. \url{https://scholarship.law.vanderbilt.edu/vlr/vol73/iss6/6}

\vspace{6pt}
Radford, A., Kim, J. W., Xu, T., Brockman, G., McLeavey, C., y Sutskever, I. (2022). Robust speech recognition via large-scale weak supervision. \textit{arXiv preprint arXiv:2212.04356}. \url{https://doi.org/10.48550/arXiv.2212.04356}

\vspace{6pt}
Radford, A., Kim, J. W., Xu, T., Brockman, G., y Amodei, D. (2023). \textit{Robust speech recognition via large-scale weak supervision}. OpenAI. \url{https://cdn.openai.com/papers/whisper.pdf}

\vspace{6pt}
Rubio Portilla, J. (2023). Audiencias virtuales y su impacto en el debido proceso penal en un distrito judicial. \textit{SCIÉNDO, 26}(4), 204--214. \url{https://doi.org/10.17268/sciendo.2023.057}

\vspace{6pt}
Russell, S. J., y Norvig, P. (2024). \textit{Inteligencia artificial: Un enfoque moderno} (4.\textordfeminine{} ed.). Pearson Educación.


\vspace{6pt}
Vaswani, A., Shazeer, N., Parmar, N., Uszkoreit, J., Jones, L., Gomez, A. N., y Polosukhin, I. (2017). Attention is all you need. \textit{Advances in Neural Information Processing Systems (NeurIPS), 30}, 5998--6008. \url{https://doi.org/10.48550/arXiv.1706.03762}

\vspace{6pt}
Vera, L. (2021). STT: Un sistema de apoyo a la transcripción de audiencias fiscales usando Vosk. \textit{Revista de Investigación de Sistemas e Informática, 24}(1), 45--57. \url{https://doi.org/10.15381/risi.v14i1.21864}

\vspace{6pt}
Yépez-Reyes, M., y Cruz-Silva, A. (2024). Inteligencia artificial en la transcripción de entrevistas: Comparativa de herramientas y métricas de calidad. \textit{Contratexto}, (41). \url{https://doi.org/10.26439/contratexto2024.n41.6750}

\end{document}
